\documentclass[11pt,a4paper]{article}

% HISTORICAL-NON-NORMATIVE: retained only for provenance; do not cite this
% narrative for any current formula, proof obligation, or theorem statement.
% CANONICAL-SOURCE: ../research/formal-theorems-and-proofs.md
% CURRENT-GENERATED-PROJECTION: theorem-lemma-proof-en.tex

\usepackage[margin=2.5cm]{geometry}
\usepackage{amsmath,amssymb,amsthm,mathtools}
\usepackage{enumitem}
\usepackage{booktabs}
\usepackage{longtable}
\usepackage[hidelinks]{hyperref}
\usepackage[T5]{fontenc}
\usepackage[utf8]{inputenc}
\usepackage[vietnamese]{babel}
\usepackage{lmodern}

% PROOF-CONTRACT: PC-2026-07-22.3; A1-A9; M1,M2,M3,M3a,M4,M5; PC-T0-PC-T6; SF-01-SF-29
% REGISTRY-ASSUMPTIONS: A1 A2 A3 A4 A5 A6 A7 A8 A9
% REGISTRY-SCOPE-LEMMAS: M1 M2 M3 M3a M4 M5
% PROOF-REGISTRY-SHA256: 2b07b56376633f6f441ef81da73ec87a4a97169a7f3b83af77bbf5f9396cd074

\newtheorem{theorem}{Định lý}
\newtheorem{lemma}{Bổ đề}
\newtheorem{assumption}{Giả thiết}
\theoremstyle{definition}
\newtheorem{definition}{Định nghĩa}
\newtheorem{remark}{Nhận xét}

\newcommand{\OCLval}{\mathrm{OCL}_{val}}
\newcommand{\ASTval}{\mathrm{AST}_{val}}
\newcommand{\BoundOCLval}{\mathrm{BoundOCL}_{val}}
\newcommand{\Obj}{\mathrm{Obj}}
\newcommand{\Class}{\mathrm{Class}}
\newcommand{\Node}{\mathrm{Node}}
\newcommand{\Rel}{\mathrm{Rel}}
\newcommand{\id}{\mathrm{id}}
\newcommand{\node}{\mathrm{node}}
\newcommand{\classNode}{\mathrm{classNode}}
\newcommand{\RepG}{\mathrm{Rep}_{G}}
\newcommand{\Adm}{\mathrm{Adm}}
\newcommand{\prop}{\mathrm{prop}}
\newcommand{\src}{\mathrm{src}}
\newcommand{\trg}{\mathrm{trg}}
\newcommand{\PhiEnc}{\Phi}
\newcommand{\boolval}{\mathrm{bool}_{val}}
\newcommand{\finiteSet}{\mathrm{finiteSet}}
\newcommand{\asSources}{\mathrm{asSources}}
\newcommand{\eqval}{\mathrm{eq}_{val}}
\newcommand{\ObjNodeVal}{\mathrm{ObjNode}_{val}}
\newcommand{\ObsObj}{\mathrm{Obs}_{obj}}
\newcommand{\ObsGraph}{\mathrm{Obs}_{graph}}
\newcommand{\liftNode}{\mathrm{lift}_{node}}
\newcommand{\encodeValue}{\mathrm{encodeValue}}
\newcommand{\returnedIds}{\mathrm{returnedIds}}
\newcommand{\Viol}{\mathrm{Viol}}
\newcommand{\sem}[2]{\llbracket #1 \rrbracket_{#2}}
\newcommand{\EvalCypher}{\mathrm{Eval}_{Cypher}}
\newcommand{\DenoteGraph}{\mathrm{Denote}_{graph}}
\newcommand{\Params}{\mathrm{Params}}
\newcommand{\ParamEnv}{\mathrm{ParamEnv}}
\newcommand{\Rows}{\mathrm{Rows}}
\newcommand{\RowCorr}{\mathrm{RowCorr}}
\newcommand{\projectSet}{\mathrm{projectSet}}
\newcommand{\reprC}{\mathrm{repr}_{C}}
\newcommand{\absC}{\mathrm{abs}_{C}}

\title{Báo cáo chứng minh chi tiết về bảo toàn ngữ nghĩa kiểm tra OCL-to-Cypher có điều kiện}
\author{}
\date{}

\begin{document}
\maketitle

\noindent\fbox{\parbox{0.94\linewidth}{\textbf{Historical draft -- not
normative.} This file predates current corrections to the certified grammar,
BR1--BR10, production-plan adequacy, and Theorems T0--T6. Use
\texttt{../research/formal-theorems-and-proofs.md} or the generated English
report for every current formula and claim.}}
\medskip

\begin{abstract}
Báo cáo này trình bày cơ sở hình thức đầy đủ cho phép chuyển invariant từ mảnh
$\OCLval$ sang truy vấn Cypher trên property graph. Kết quả trung tâm là tương
đương có điều kiện giữa tập đối tượng UML vi phạm invariant và tập đối tượng có
định danh được truy vấn sinh ra trả về. Chứng minh phân rã chuỗi biến đổi thành
các nghĩa vụ về biểu diễn Source--AST, binding và định kiểu, trừu tượng hóa
Validation Algebra, chuẩn hóa, tương ứng mô hình đối tượng--đồ thị và hiện thực
hóa Cypher. Báo cáo phát biểu tường minh các giả thiết A1--A9, hợp đồng biểu
diễn R1--R7 và BR1--BR7, các tiên đề Cypher CY1--CY9, các bổ đề cục bộ cùng bảy
định lý T0--T6. Phạm vi chỉ bao phủ ngữ nghĩa kiểm tra theo tập hữu hạn và
không được hiểu là chứng minh cho toàn bộ OMG OCL, mọi phép mã hóa graph hoặc
Cypher tùy ý.
\end{abstract}

\tableofcontents
\newpage

\section{Phát biểu đóng góp và phạm vi}

Tài liệu chứng minh phát biểu có điều kiện sau:

\begin{quote}
Chúng tôi chứng minh sự bảo toàn đầu-cuối có điều kiện của tập đối tượng vi
phạm đối với một mảnh kiểm tra OCL theo ngữ nghĩa tập hữu hạn, trên một phép mã
hóa UML sang đồ thị thuộc tính được phát biểu tường minh.
\end{quote}

Kết quả được giới hạn có chủ ý. Tài liệu không tuyên bố tính đúng đắn cho toàn
bộ OMG OCL, ngữ nghĩa đầy đủ của null/invalid, Bag/Sequence/OrderedSet, tập hợp
có thứ tự, Cypher tùy ý, phép mã hóa đồ thị tùy ý hoặc các cấu trúc chỉ được xử
lý bằng cơ chế dự phòng.

\begin{assumption}[Phạm vi chứng minh]
Mọi định lý dưới đây đều giả thiết rằng:
\begin{enumerate}[label=(A\arabic*)]
  \item $e\in\OCLval$, $ast=\mathrm{Parse}(e)$ thành công và
  $decode_{val}(ast)=e$, nên $ast\in\ASTval$.
  \item $b=T_{\mathrm{BIND}}(ast,MM)$ thành công và $b$ là một biểu thức
  invariant định kiểu đúng trong $\BoundOCLval(MM)$. Ở runtime, đối số này có
  thể được cung cấp bởi một view siêu mô hình phía đồ thị tương đương resolver
  với $MM$.
  \item Kiểm tra sử dụng ngữ nghĩa tập hữu hạn.
  \item $MM=M_2$ là siêu mô hình miền hữu hạn và $M=M_1$ là mô hình đối tượng
  UML hữu hạn, hợp lệ, phù hợp với $MM$.
  \item $G=\PhiEnc(MM,M)$ thỏa hợp đồng mã hóa đồ thị repository: $G$ lưu cả
  phần siêu mô hình $G_{M_2}$, phần instance $G_{M_1}$ và các cạnh typing giữa
  chúng; view siêu mô hình đọc từ $G_{M_2}$ đầy đủ và tương đương resolver với
  $MM$.
  \item $T_{\mathrm{TEXT}}$ chỉ sinh mảnh Cypher được hỗ trợ.
  \item Truy vấn Cypher được sinh là truy vấn kiểm tra invariant, trả về các đối
  tượng vi phạm bằng định danh ổn định.
  \item $ScalarClosed(e,M)$ đúng với mọi phép đánh giá con vô hướng có thể đạt tới.
  \item $BottomSeparated(MM,M,e)$ đúng: biểu diễn dành riêng cho bottom không
  giao với mọi giá trị mô hình, định danh, khóa siêu mô hình và payload
  qualifier có thể đạt tới.
\end{enumerate}

Trong adapter chuẩn, A9 được giải phóng theo construction: scalar và khóa mô
hình được chấp nhận chỉ thuộc Boolean, Int64, Real64 hữu hạn hoặc String, entity
được biểu diễn bằng node, còn bottom trong tập là map tag riêng
\texttt{\{\_\_oclBottom:true\}}. Kho thuộc tính và codec qualifier được chứng
nhận không nhận map làm scalar mô hình. Verifier fixture kiểm tra điều kiện miền
này trước khi coi A9 đã được giải phóng. A8 vẫn là premise theo từng
model/expression vì số học động có thể tràn hoặc sinh mẫu chia không xác định dù
các toán hạng đều định kiểu đúng.
\end{assumption}

\subsection*{Hợp đồng định lý chuẩn}

Phiên bản hợp đồng là \texttt{PC-2026-07-22.3}. Các dòng dưới đây là hình chiếu
LaTeX đã đồng bộ từ canonical JSON block trong
\texttt{md/research/formal-theorems-and-proofs.md}; phần
diễn giải của định lý có thể mở rộng một dòng nhưng không được loại bỏ giả
thiết hoặc phụ thuộc của dòng đó.

% BEGIN GENERATED PROOF-CONTRACT-KERNEL
\begin{longtable}{p{1.2cm}p{4.5cm}p{2.4cm}p{5.0cm}}
\toprule
ID & Statement kernel & Assumptions & Required results\\
\midrule
PC-T0 & Các quan sát kiểm tra trên đối tượng và đồ thị đủ điều kiện trùng khớp & A4, A5 & R1, R2, R3, R4, R5, R6, R7\\
PC-T1 & Ngữ nghĩa Source đã giải mã và Bound duy nhất tương ứng sau khi xóa thẻ & A1, A2, A4, A8 & M2, M3, M3a, M4, B0, B1, B2, B3, B4, SB0\\
PC-T2 & Ngữ nghĩa Bound sau xóa thẻ bằng ngữ nghĩa VA có kiểu phía đối tượng & A2, A3, A4, A8 & M2, M3, M3a, VA1, VA2, VA3, BV0, B4\\
PC-T3 & Chuẩn hóa bottom-up dừng, bảo toàn kiểu/ngữ nghĩa/miền đóng và đạt NF\_R theo alpha & A2, A3, A8 & N0, N1, N2, N3, N4\\
PC-T4 & Đánh giá VA có kiểu giao hoán với mã hóa trên các đánh giá reachable đóng & A2, A3, A4, A5, A8 & T0, G1, G2, G3, G3a, G4\\
PC-T5 & Truy vấn tham chiếu NVA--CQM--Cypher trả về chính xác ID vi phạm VA đồ thị và không có ID ma & A2, A3, A5, A6, A7, A8, A9 & C1, C2, C3, C3a, C4, C5, C5a, C5b, C6, SpecPlanAdequacy, SpecPlanSim_sound, TXT1, TXT2, TXT3, TXT4, TXT5, BR1-BR10, CY1-CY9\\
PC-T6 & Tập ID trả về bằng ảnh ID của tập vi phạm OCL & A1, A2, A3, A4, A5, A6, A7, A8, A9 & T0, T1, T2, T3, T4, T5\\
\bottomrule
\end{longtable}
% PROTOTYPE-CORRECTNESS-CLAIM: CONDITIONAL
% BLOCKING-OPEN-OR-PARTIAL: PO-18
% RUNTIME-PROFILE: RP-NEO4J-2026.06-CYPHER5-CANONICAL-V1
% END GENERATED PROOF-CONTRACT-KERNEL

\section{Siêu mô hình nguồn $\OCLval$, biểu diễn AST và $\BoundOCLval$}

$\OCLval$ là siêu mô hình của mảnh ngôn ngữ nguồn và được định nghĩa trước biểu
diễn parser. Nó quy định các constructor nguồn được hỗ trợ nhưng chưa chứa tham
chiếu UML đã phân giải. Hàm $encode_{ast}:\OCLval\to\ASTval$ tạo biểu diễn AST,
còn hàm bộ phận $decode_{val}:AST_{OCL}\rightharpoonup\OCLval$ nhận ra đúng các
AST thuộc mảnh nguồn. Sau đó, và chỉ sau binding thành công, ta thu được mô hình
đã phân giải và định kiểu $\BoundOCLval(MM)$. Vì vậy thứ tự đặc tả là
$\OCLval\to\ASTval\to\BoundOCLval(MM)$, trong khi pipeline thực thi vẫn là
OCL Text $\to$ Parse $\to$ OCL AST $\to T_{\mathrm{BIND}}$ Bound OCL.

\subsection{Siêu mô hình nguồn $\OCLval$}

Siêu mô hình nguồn chứa tên cú pháp chưa phân giải. Với tên biến $x$, tên lớp
$cn$, tên property $pn$ và danh sách qualifier nguồn $\bar q$, các biểu thức
nguồn được sinh bởi:
\[
\begin{array}{rcl}
s &::=& self \mid x \mid c \mid Set\{s_1,\ldots,s_n\} \mid s.pn[\bar q]
 \mid cn.allInstances() \mid not\ s \mid s\ op\ s\\
&\mid& if\ s\ then\ s\ else\ s
 \mid let\ x:T=s\ in\ s
 \mid s\to iop(x\mid s)\\
&\mid& s\to cop(s?)
 \mid s.oclIsKindOf(cn)
 \mid s.oclAsType(cn).
\end{array}
\]
với $n\geq 1$ trong trường hợp Set literal.
Ở tầng này, $pn$ chưa được phân loại thành attribute hay association role,
$cn$ chưa phải tham chiếu lớp và navigation chưa có association, role hoặc
direction. Các thông tin đó chỉ xuất hiện trong $\BoundOCLval(MM)$.

\subsection{Các lớp siêu mô hình và đặc trưng cấu trúc}

Các lớp siêu mô hình trong bảng sau thuộc $\BoundOCLval(MM)$. Mọi tham chiếu có
đối miền thuộc $MM$ đều là tham chiếu đến phần tử mô hình đã
phân giải, không phải tên String không định danh đầy đủ. Mọi containment đều là
hợp thành và mỗi biểu thức có đúng một container; vị từ gốc nằm trong invariant,
còn invariant nằm trong module.

\begingroup\small
\begin{longtable}{p{0.25\textwidth}p{0.47\textwidth}p{0.22\textwidth}}
\toprule
Lớp siêu mô hình & Đặc trưng cấu trúc & Lực lượng/kiểu\\
\midrule
$BoundValModule$ & $invariants$ & containment $0..*$ của $BoundValInvariant$\\
$BoundValInvariant$ & $contextClass$, $name$, $predicate$ & $Class(MM)$ đã phân giải, String, containment $1$\\
$BoundValExpr\langle\tau\rangle$ (trừu tượng) & $staticType$ & đúng một $\tau$\\
$VSelf$ & không có đặc trưng bổ sung & $ValExpr\langle C\rangle$\\
$VVariable$ & $declaration$ & một khai báo từ vựng có kiểu $\tau$\\
$VLiteral$ & $kind$, $value$ & một literal nguyên thủy được hỗ trợ\\
$VSetLiteral$ & $elements$ & containment $1..*$ của các phần tử cùng kiểu chuẩn\\
$VAttribute$ & $source$, $attribute$ & containment $1$; $Attribute(C,\tau)$ đã phân giải\\
$VNavigation$ & $source$, $association$, $fromRole$, $toRole$, $direction$, $qualifiers$ & containment $1$ và $0..*$; metadata nhị phân đã phân giải\\
$VAllInstances$ & $referredClass$ & một tham chiếu $Class(MM)$\\
$VNot$ & $operand$ & containment $1$\\
$VBinary$ & $operator$, $left$, $right$ & toán tử và hai containment\\
$VIf$ & $condition$, $thenBranch$, $elseBranch$ & ba containment\\
$VLet$ & $declaration$, $init$, $body$ & khai báo và hai containment\\
$VIterator$ & $operation$, $declaration$, $source$, $body$ & thẻ iterator và hai containment\\
$VCollectionOp$ & $operation$, $source$, $argument?$ & source và nhiều nhất một đối số\\
$VTypeOp$ & $operation$, $source$, $referredClass$ & chỉ kind-of hoặc cast\\
\bottomrule
\end{longtable}
\endgroup

Các miền liệt kê đóng là:
\[
\begin{split}
BinaryOp={}&\{and,or,xor,implies,=,<>,<,\leq,>,\geq,+,-,*,/\},\\
IteratorOp={}&\{exists,forAll,select,reject,collect,isUnique\},\\
CollectionOp={}&\{includes,excludes,includesAll,excludesAll,\\
&\hspace{1.4cm}union,intersection,asSet,size,isEmpty,notEmpty\},\\
TypeOp={}&\{kindOf,cast\},\qquad Direction=\{forward,reverse\}.
\end{split}
\]
Đặc biệt, $typeOf$ không phải phần tử của $TypeOp$.

\subsection{Cú pháp trừu tượng của $\BoundOCLval(MM)$}

Với $C,D\in Class(MM)$ và biến có kiểu $x$, thành phần biểu thức là họ nhỏ nhất
được sinh bởi:
\[
\begin{array}{rcl}
e &::=& self \mid x \mid c \mid SetLiteral([e_1,\ldots,e_n]) \mid Attribute(e,a)\\
  &\mid& Navigation(e,A,fr,tr,dir,[q_1,\ldots,q_k])\\
  &\mid& AllInstances(C) \mid Not(e) \mid Binary(op,e,e)\\
  &\mid& If(e,e,e) \mid Let(x:T,e,e)\\
  &\mid& Iterator(iop,e,x,e) \mid CollectionOp(cop,e,e?)\\
  &\mid& TypeKindOf(e,C) \mid Cast(e,C).
\end{array}
\]
với $n\geq 1$ trong trường hợp $SetLiteral$.
Gốc duy nhất được Định lý~6 bao phủ là $Invariant(C,R,p)$ với $p:Boolean$.
Truy vấn mức cao nhất của prototype không phải gốc invariant của $\OCLval$.

\subsection{Các ràng buộc hợp lệ cấu trúc}

Một thể hiện thuộc $\BoundOCLval(MM)$ khi và chỉ khi mọi điều kiện sau đều đúng.

\begin{enumerate}[label=(MV-WF\arabic*)]
  \item Đồ thị containment hợp thành là hữu hạn, có gốc, không chu trình và
  không chia sẻ nút biểu thức.
  \item Mỗi invariant có lớp ngữ cảnh đã phân giải và vị từ Boolean được định
  kiểu trong môi trường chứa $self:C$.
  \item Mỗi biểu thức có kiểu chuẩn duy nhất do các luật định kiểu gán; $\bot$
  ở thời gian chạy là một giá trị và không bao giờ là kiểu tĩnh.
  \item Tham chiếu thuộc tính và lớp thuộc $MM$. Mỗi navigation là navigation
  liên kết nhị phân được phân giải duy nhất, với role, hướng, lớp đích, số lượng
  qualifier và kiểu qualifier cố định.
  \item Biểu thức có giá trị collection chỉ dùng $Set(\tau)$. Navigation được
  nâng thống nhất thành $Set(C)$, kể cả đầu mút UML có bội số $0..1$ hoặc $1$.
  \item Khai báo iterator và let chỉ có phạm vi trong thân của nó. Tra cứu chọn
  khai báo gần nhất và đổi tên alpha tránh bắt giữ biến.
  \item Thân $collect$ có kiểu không phải collection. Kết quả là ảnh của tập
  hữu hạn; không biểu diễn flatten ngầm, bội số hoặc thứ tự.
  \item Qualifier là danh sách có thứ tự, lực lượng cố định của các biểu thức vô
  hướng nguyên thủy được hỗ trợ, với phép tuần tự hóa đơn ánh do profile mã hóa
  quy định.
  \item Set literal là hữu hạn, không rỗng, không lồng collection và có một kiểu
  phần tử chuẩn suy ra duy nhất. Thân $isUnique$ không có kiểu collection;
  $union$, $intersection$ và $asSet$ chỉ thao tác trên $Set$ tương thích.
  \item Chỉ xuất hiện các miền toán tử đóng ở trên. Do đó $any$, $one$,
  $sortedBy$, $iterate$, $closure$, $count(element)$,
  \texttt{oclIsTypeOf} và ràng buộc không phải invariant không có trường hợp
  trong siêu mô hình.
\end{enumerate}

\begin{lemma}[M1 Biểu diễn Source--AST]
Hai hàm $encode_{ast}$ và $decode_{val}$ tạo thành một retraction:
\[
decode_{val}(encode_{ast}(s))=s\qquad(s\in\OCLval).
\]
Mỗi $ast\in\ASTval$ có đúng một biểu thức nguồn được decode, tới phép đổi tên
alpha tránh bắt giữ biến bị ràng buộc.
\end{lemma}
\begin{proof}
Định nghĩa $encode_{ast}$ bằng đệ quy cấu trúc: mỗi constructor nguồn ánh xạ
sang đúng một constructor AST và mã hóa đệ quy các phần tử con. Ngược lại,
$decode_{val}$ có đúng một phương trình cho mỗi dạng AST được hỗ trợ và không
xác định trên mọi dạng khác. Quy nạp trên cây nguồn hữu hạn chứng minh phương
trình retraction; quy nạp trên AST được chấp nhận chứng minh tính duy nhất. Tên
binder chỉ có thể khác bởi đổi tên alpha tránh bắt giữ.
\end{proof}

\begin{lemma}[M2 Tính đóng của biểu thức Bound con và quy nạp cấu trúc]
Mọi biểu thức được chứa trong một biểu thức $\BoundOCLval(MM)$ tự nó là biểu thức
$\BoundOCLval(MM)$ hợp lệ và định kiểu đúng dưới môi trường từ vựng do các nút tổ
tiên tạo ra. Hơn nữa, quan hệ containment nghiêm ngặt là chính quy tốt.
\end{lemma}
\begin{proof}
MV-WF1 làm containment nghiêm ngặt thành quan hệ hữu hạn không chu trình, do đó
là quan hệ chính quy tốt. Ràng buộc định kiểu của lớp siêu mô hình yêu cầu mọi
tiền đề của nút con trước khi chấp nhận nút cha. Với thân iterator và let,
MV-WF6 thêm đúng binding đã khai báo; các nút con khác giữ môi trường của nút
cha. Quy nạp theo đường đi duy nhất từ gốc chứng minh tính đóng và môi trường
phù hợp cho từng nút con. Vì thế quy nạp cấu trúc và đệ quy cấu trúc toàn phần
là hợp lệ trên các thể hiện $\BoundOCLval$.
\end{proof}

\begin{lemma}[M3 Tính duy nhất của kiểu chuẩn]
Nếu biểu thức được chấp nhận $e$ có các suy dẫn định kiểu chuẩn
$MM;\Gamma\vdash e:\tau_1$ và $MM;\Gamma\vdash e:\tau_2$, thì
$\tau_1=\tau_2$.
\end{lemma}
\begin{proof}
Quy nạp trên lớp siêu mô hình của $e$. Kiểu cơ sở được xác định bởi luật
literal, khai báo hoặc lớp ngữ cảnh. Kết quả thuộc tính, navigation,
all-instances và toán tử kiểu được cố định bởi tham chiếu đã phân giải có tính
hàm. Toán tử Boolean và collection có kiểu kết quả cố định. Số học dùng hàm bộ
phận $arithResult$, còn định kiểu nhánh dùng hàm bộ phận $join$; cả hai đều có
tính hàm. Set literal và hợp/giao dùng hàm bộ phận $setJoin$ có tính hàm theo
M3a; $asSet$ giữ kiểu nguồn và $isUnique$ trả Boolean. Trường hợp iterator và let dùng giả thiết quy nạp dưới phép mở rộng
môi trường duy nhất từ MV-WF6. Không có luật subsumption tổng quát làm thay đổi
chú thích kiểu chuẩn. Do đó hai kiểu kết quả bằng nhau.
\end{proof}

\begin{lemma}[M4 Tính đúng đắn của admission và binding]
Nếu $decode_{val}(ast)=s$ và binder được chứng nhận suy dẫn
$MM;\Gamma\vdash ast\Rightarrow b:\tau$, thì $s\in\OCLval$,
$ast\in\ASTval$, $b\in\BoundOCLval(MM)$ và $MM;\Gamma\vdash b:\tau$.
\end{lemma}
\begin{proof}
Quy nạp trên suy dẫn binding thành công. Mỗi luật binder tạo đúng một lớp siêu
mô hình đã liệt kê. Giả thiết quy nạp cho nút con thiết lập containment hữu hạn
có kiểu. Resolver có tính hàm thiết lập tham chiếu đã phân giải và metadata
navigation duy nhất; luật iterator và let dùng phép mở rộng từ vựng đã nêu;
luật collect và isUnique buộc thân không có kiểu collection. B-SetLit và
B-SetComb chỉ áp dụng khi $setJoin$ xác định duy nhất, còn B-AsSet giữ kiểu
phần tử. Kiểu kết quả của luật là kiểu
của lớp siêu mô hình tương ứng. Vì bảng được chứng nhận không có luật cho toán
tử bị loại, MV-WF10 được bảo toàn. Các trường hợp là đầy đủ vì phán đoán binding
được định hướng cú pháp.
\end{proof}

\begin{lemma}[M5 Tính ổn định của phạm vi loại trừ]
Không cặp source/AST được chấp nhận nào và không kết quả
$\BoundOCLval(MM)$ tương ứng nào có thể chứa toán tử bị loại, kiểu collection
nhạy với thứ tự hoặc bội số, navigation không nhị phân hoặc gốc không phải invariant.
\end{lemma}
\begin{proof}
Theo M1, không có phương trình $decode_{val}$ cho constructor AST bị loại. Theo
M2, mỗi biểu thức Bound con có một lớp siêu mô hình đã liệt kê. Miền toán tử
đóng, kiểu collection chỉ do $Set$ sinh và metadata navigation là nhị phân.
Vì vậy không có trường hợp biểu diễn nguồn, AST hoặc Bound cho construct bị loại.
\end{proof}

M1--M5, bao gồm M3a, cung cấp chứng nhận phạm vi cho các chứng minh phía sau. M2 biện minh cho
quy nạp cấu trúc Bound, M3 cung cấp tính hàm của kiểu, M4 nối admission AST với
$\BoundOCLval(MM)$, và F1 bên dưới chứng minh biểu diễn Source--AST cùng erase--rebind theo
chiều ngược dưới các resolver ổn định. Chúng không thêm constructor hoặc mở
rộng Định lý~6.

\begin{definition}[Chân lý kiểm tra]
Vị từ chân lý ở mức kiểm tra được định nghĩa bởi:
\[
\boolval(v)=\mathrm{true} \iff v=\mathrm{true}.
\]
Mọi giá trị khác, gồm bottom, giá trị giống null, giống invalid hoặc giá trị
runtime không được hỗ trợ, đều không được xem là đúng theo kiểm tra. Đây là một
chính sách kiểm tra, không phải hình thức hóa đầy đủ các bảng chân lý OMG OCL.
\end{definition}

\begin{definition}[Phép hoàn chỉnh toàn phần ở mức kiểm tra]
Trong miền định lý, hàm ngữ nghĩa là toàn phần vào vũ trụ giá trị có gắn thẻ.
Truy cập thuộc tính trên $\bot$ cho $\bot$; phép nâng nguồn navigation ánh xạ
$\bot$ thành tập rỗng; phép toán số học không xác định, kể cả chia cho không,
cho $\bot$; phép so sánh thứ tự không xác định cũng cho $\bot$. Đẳng thức có
kiểu là toàn phần:
\[
\eqval(\bot,\bot)=true,
\qquad
\eqval(\bot,v)=false\quad(v\neq\bot).
\]
Các vị từ kiểu trên $\bot$ cho false và ép kiểu $\bot$ cho $\bot$.
$Collect$ giữ $\bot$ như một phần tử ảnh. Kiểu thân được chấp nhận của nó không
phải collection, vì vậy không tuyên bố phép làm phẳng ngầm hay hiện thực hóa
collection lồng nhau. Các mệnh đề này là chính sách kiểm tra, không phải ngữ
nghĩa invalid của OMG OCL.
\end{definition}

\begin{definition}[Nâng nguồn navigation và đẳng thức có kiểu]
Receiver của navigation sử dụng
\[
\asSources(\mathrm{Set}(S))=S,\qquad
\asSources(v)=\{v\}\ \text{khi }v\text{ là một Entity},\qquad
\asSources(\bot)=\emptyset.
\]
Định kiểu navigation nhị phân bảo đảm mọi phần tử được nâng đều là Entity.
Các miền giá trị là những hợp rời nhau có gắn thẻ. Đẳng thức có kiểu
$\eqval(v_1,v_2)$ đúng khi và chỉ khi hai toán hạng có cùng loại giá trị và biểu
thị cùng giá trị; đẳng thức thực thể là định danh đối tượng ở phía đối tượng và
định danh nút ở phía đồ thị. So sánh thứ tự vẫn chỉ áp dụng cho các kiểu vô
hướng tương thích.
\end{definition}

\begin{definition}[Hàm hỗ trợ tập hữu hạn]
Các toán tử collection sử dụng:
\[
\finiteSet(\mathrm{Set}(S))=S,\qquad
\finiteSet(\bot)=\emptyset.
\]
Với các biểu thức thuộc phạm vi định lý, ngữ nghĩa có giá trị collection là tập
hữu hạn hoặc bottom. Việc xem bottom như tập rỗng tại vị trí collection là một
phần của chính sách ở mức kiểm tra.
\end{definition}

\begin{definition}[Giá trị và môi trường có kiểu]
Với $I\in\{obj,graph\}$, các giá trị được chỉ số hóa theo kiểu tĩnh:
\[
Val_I(Boolean)=\{true,false\},\quad
Val_I(Integer),\ Val_I(Real),\ Val_I(String),
\]
\[
Val_{obj}(C)=\{o\mid classOf(o)\mathrel{conformsTo}C\},\qquad
Val_{graph}(C)=\{\node(o)\mid o\in Val_{obj}(C)\},
\]
\[
Val_I(Set(\tau))=\mathcal P_{fin}(Val_I(\tau)\cup\{\bot\}),\qquad
Val_I^\bot(\tau)=Val_I(\tau)\cup\{\bot\}.
\]
Ta viết $\rho\models_M\Gamma$ khi và chỉ khi $\rho$ và $\Gamma$ có cùng miền,
đồng thời $\rho(x)\in Val_{obj}^\bot(\Gamma(x))$ với mọi $x$; định nghĩa
$\eta\models_G\Gamma$ tương tự. Ngoài ra,
\[
\rho\sim_\Phi\eta\iff
\operatorname{dom}(\rho)=\operatorname{dom}(\eta)\ \land\
\forall x:\eta(x)=\encodeValue(\rho(x)).
\]
Profile vô hướng chính xác là
\[
Int64=\{n\in\mathbb Z\mid -2^{63}\le n\le 2^{63}-1\},
\]
và $Real64$ là tập các giá trị IEEE-754 binary64 hữu hạn, trong đó đẳng thức số
đồng nhất hai số không có dấu khác nhau. Phép toán số thực làm tròn tới giá trị
gần nhất, trường hợp hòa làm tròn về số chẵn, và chỉ được chấp nhận khi kết quả
sau làm tròn là hữu hạn. Các phép $+,-,*$ trên số nguyên là phép toán toán học
có kiểm tra trong $Int64$; phép chia yêu cầu mẫu khác không và trả về $Real64$.
Ép số nguyên sang số thực chỉ được chấp nhận với $|n|\le 2^{53}$, nên phép
chuyển đổi là chính xác. Chuỗi là dãy hữu hạn các scalar value Unicode, có đẳng
thức chính xác và không có thứ tự. Đẳng thức thực thể là đẳng thức định danh.
$ScalarClosed(e,M)$ nghĩa là mọi đánh giá con vô hướng có thể đạt tới, kể cả
các mở rộng iterator và let, đều thỏa các tiền đề này. Tràn số, NaN, vô cực,
chia cho không, ép kiểu không chính xác, đối chiếu phụ thuộc locale và các loại
vô hướng không được hỗ trợ nằm ngoài định lý.
\end{definition}

\begin{definition}[Chính sách thống nhất cho kết quả navigation]
Mọi navigation nhị phân được chấp nhận đều có kiểu tĩnh $Set(C)$. Một
navigation UML có bội số nguồn $0..1$ hoặc $1$ được nâng thành tập rỗng hoặc tập
đơn. Đây là chính sách kiểm tra theo tập hữu hạn của $\OCLval$, không phải định
kiểu kết quả collection đầy đủ của OMG OCL.
\end{definition}

\section{Ký hiệu cốt lõi}

\begin{longtable}{ll}
\toprule
Ký hiệu & Ý nghĩa\\
\midrule
$MM$ & siêu mô hình nguồn dạng UML\\
$M$ & mô hình đối tượng phù hợp với $MM$\\
$G$ & đồ thị thuộc tính\\
$\PhiEnc(MM,M)$ & phép mã hóa đồ thị của $MM$ và $M$\\
$\Obj(M)$ & tập hữu hạn các đối tượng trong $M$\\
$\Class(MM)$ & tập hữu hạn các lớp trong $MM$\\
$\Node(G)$, $\Rel(G)$ & các nút và quan hệ của đồ thị\\
$\id(o)$ & định danh đối tượng ổn định\\
$\node(o)$ & nút đồ thị biểu diễn đối tượng $o$\\
$\classNode(C)$ & nút đồ thị biểu diễn lớp $C$\\
$\RepG(o,n)$ & nút $n$ được $\PhiEnc$ sinh để biểu diễn đối tượng $o$\\
$\ObjNodeVal(G)$ & các nút đối tượng đồ thị nhìn thấy bởi kiểm tra\\
\bottomrule
\end{longtable}

\subsection*{Danh mục hàm ngữ nghĩa chuẩn}

Đây là bảng chữ ký chuẩn tắc duy nhất trong tài liệu. Các định danh của bảng
được đồng bộ với phiên bản registry \texttt{PC-2026-07-22.3}. Các lần sử dụng
phía sau chỉ được phép bỏ chỉ số trên về kiểu khi suy dẫn định kiểu xác định duy
nhất chỉ số đó; việc lược bỏ này không tạo ra một hàm ngữ nghĩa thứ hai.

\scriptsize
\begin{longtable}{p{1.15cm}p{3.0cm}p{8.8cm}}
\toprule
ID & Ký hiệu & Miền và đối miền\\
\midrule
SF-01 & $\sem{\cdot}{S}$ & $Expr_S(\tau)\times Model\times Env_S\to Val_S^\bot(\tau)$\\
SF-02 & $\sem{\cdot}{B}$ & $Expr_B(\tau)\times Model\times Env_B\to Val_B^\bot(\tau)$\\
SF-03 & $erase_S$ & $Val_S^\bot(\tau)\to Val_{obj}^\bot(\tau)$\\
SF-04 & $erase_B$ & $Val_B^\bot(\tau)\to Val_{obj}^\bot(\tau)$\\
SF-05 & $\sem{\cdot}{OCLval}$ & $Expr_S(\tau)\times Model\times Env_S\to Val_{obj}^\bot(\tau)$, được định nghĩa bởi $erase_S\circ\sem{\cdot}{S}$\\
SF-06 & $\sem{\cdot}{Bound}$ & $Expr_B(\tau)\times Model\times Env_B\to Val_{obj}^\bot(\tau)$, được định nghĩa bởi $erase_B\circ\sem{\cdot}{B}$\\
SF-07 & $\sem{\cdot}{obj}^{\tau}$ & $Expr_{VA}(\tau)\times Model\times Env_{obj}\to Val_{obj}^\bot(\tau)$\\
SF-08 & $\sem{\cdot}{graph}^{\tau}$ & $Expr_{VA}(\tau)\times Graph\times Env_{graph}\to Val_G^\bot(\tau)$\\
SF-09 & $\DenoteGraph^{\tau}$ & Bí danh của SF-08 giới hạn cho các accessor đồ thị nguyên thủy\\
SF-10 & $\encodeValue$ & $Val_{obj}^\bot(\tau)\to Val_G^\bot(\tau)$\\
SF-11 & $\finiteSet_\tau$ & $Val_I^\bot(Set(\tau))\to\mathcal P_{fin}(Val_I^\bot(\tau))$\\
SF-12 & $\asSources_C$ & receiver navigation dạng Entity/tập/bottom $\to\mathcal P_{fin}(Val_I(C))$\\
SF-13 & $\boolval$ & $Val_I^\bot(Boolean)\to Boolean$\\
SF-14 & $\eqval^\tau$ & $Val_I^\bot(\tau)^2\to Boolean$\\
SF-15 & $I_{attr}^{\tau}$ & $Val_I^\bot(C)\times Attribute(C,\tau)\to Val_I^\bot(\tau)$\\
SF-16 & $I_{nav}$ & receiver có kiểu, metadata liên kết và giá trị qualifier $\to\mathcal P_{fin}(Val_I(D))$\\
SF-17 & $I_{class}$ & $Class\to\mathcal P_{fin}(Val_I(C))$\\
SF-18 & $\reprC$ & $Val_G^C\to CyVal_C$\\
SF-19 & $\absC$ & $CyVal_C\to Val_G^C$, với $\absC(\reprC(v))=v$\\
SF-20 & $EvalExpr_{G,\sigma}$ & $CExpr\times Row\to CyVal_C$\\
SF-21 & $EvalRows_{G,\sigma}$ & $CQuery\times Bag(Row)\to Bag(Row)$\\
SF-22 & $\EvalCypher$ & $CQuery\times Graph\times ParamEnv\to Bag(Row)$\\
SF-23 & $ExecPrelude$ & $List(Clause)\times Row\to Bag(Row)$\\
SF-24 & $\projectSet$ & $Alias\times Bag(Row)\to\mathcal P_{fin}(Val_G^C)$\\
SF-25 & $params$ & $CQuery\to ParamEnv$\\
SF-26 & $\Params$ & $CQuery\to\mathcal P(ParamName)$\\
SF-27 & $\returnedIds$ & $CQuery\times Graph\times ParamEnv\to\mathcal P(ObjectId)$\\
SF-28 & $\Viol$ & $Expr_S(Boolean)\times Class\times Model\to\mathcal P(Obj(M))$\\
SF-29 & $navCollectionView$ & $Val_I^\bot(C)\to\mathcal P_{fin}(Val_I(C))$, singleton/empty view chỉ tại biên collection trực tiếp của navigation to-one\\
\bottomrule
\end{longtable}
\normalsize

Với truy vấn được sinh, dạng viết tắt hai đối số là
$\returnedIds(q,G):=\returnedIds(q,G,params(q))$.

Trong toàn bộ chứng minh, $MM=M_2$ ký hiệu siêu mô hình miền và $M=M_1$ ký hiệu
mô hình instance phù hợp. $\PhiEnc$ ký hiệu hàm mã hóa repository và
$G=\PhiEnc(MM,M)$ ký hiệu đồ thị cụ thể đã mã hóa. Về mặt quan sát, ta phân rã
\[
G=G_{M_2}\cup G_{M_1}\cup G_{typing},
\]
trong đó $G_{M_2}$ lưu khai báo lớp, thuộc tính, liên kết, đầu mút, role,
qualifier và kế thừa; $G_{M_1}$ lưu đối tượng, giá trị thuộc tính và liên kết
instance; còn $G_{typing}$ nối đối tượng M1 với lớp M2. Đây là các phép chiếu
logic của cùng một repository graph, không bắt buộc là ba cơ sở dữ liệu vật lý.
Ta dùng $G$ cho đánh giá phía đồ thị để tránh ký hiệu tự tham chiếu.

Đặt $GraphMMView(G_{M_2})$ là giao diện resolver đọc metadata M2 đã mã hóa.
Hợp đồng A5 yêu cầu view này tương đương resolver với $MM$: mọi lời gọi có thể
đạt tới của $resolveClass$, $resolveAttribute$, $resolveNavigation$,
$resolveInheritance$ và $resolveQualifier$ đều trả cùng phần tử hoặc cùng trạng
thái không xác định. Vì vậy, ở runtime có thể dùng
\[
T_{\mathrm{BIND}}(ast,GraphMMView(G_{M_2}))
\]
thay cho $T_{\mathrm{BIND}}(ast,MM)$ mà không đổi kết quả Bound, tới đồng cấu
tham chiếu. Đây là đọc lại schema được mã hóa tường minh, không phải suy diễn
schema từ các link instance đang có trong $G_{M_1}$.

Chuỗi biến đổi thành công là:
\[
\mathrm{ast}=\mathrm{Parse}(e),\quad
s=decode_{val}(\mathrm{ast}),\quad
b=T_{\mathrm{BIND}}(\mathrm{ast},MM),\quad
va=T_{\mathrm{VA}}(b),\quad
nva=T_{\mathrm{NORM}}(va),\quad
q=T_{\mathrm{TEXT}}(nva).
\]

Ba miền được nối bởi các hàm sau:
\[
encode_{ast}:\OCLval\to\ASTval,\qquad
decode_{val}:AST_{OCL}\rightharpoonup\OCLval,
\]
\[
decode_{val}(encode_{ast}(s))=s,\qquad
T_{\mathrm{BIND}}:\ASTval\times MM\rightharpoonup\BoundOCLval(MM),
\]
\[
\Adm_{MM}(ast)\iff ast\in\ASTval\ \land\
T_{\mathrm{BIND}}(ast,MM)\downarrow.
\]

Hợp đồng mã hóa còn yêu cầu khóa tra cứu lớp đơn ánh và phép tuần tự hóa đơn ánh
cho danh sách qualifier nguyên thủy có thứ tự:
\[
key(C_1)=key(C_2)\iff C_1=C_2,
\qquad
encodeQualifierList(\bar q_1)=encodeQualifierList(\bar q_2)
\iff \bar q_1=\bar q_2.
\]
Phép tuần tự hóa này hoạt động theo từng thành phần, bảo toàn thứ tự cùng số
ngôi qualifier đã khai báo và chỉ áp dụng cho các loại qualifier vô hướng
nguyên thủy được định lý hỗ trợ. Mỗi nút lớp có nhãn \texttt{Class} và lưu
\texttt{classKey}$=key(C)$. Chỉ khi tên lớp là duy nhất toàn cục trong $MM$ mới
được thay khóa này bằng \texttt{name} không định tính; nếu không, các pattern
context, all-instances và kind-of phải dùng \texttt{classKey}. Kiểu chính xác
không thuộc mảnh được chứng nhận; việc bổ sung nó đòi hỏi một quan sát kiểu
runtime trực tiếp, đơn ánh và các trường hợp chứng minh riêng.
Do đó $\OCLval$ là mảnh ngôn ngữ nguồn; $\ASTval$ là biểu diễn parser; và
$\BoundOCLval(MM)$ là siêu mô hình định kiểu đã phân giải.

\begin{definition}[Tương đương quan sát kiểm tra]
Định nghĩa các bộ quan sát có kiểu
\[
\ObsObj(MM,M)=(Obj,Id,Type,Attr,Link,Nav,AllInst)_{obj},
\]
\[
\ObsGraph(MM,G)=(ObjNode,Id,Type,Attr,Link,Nav,AllInst)_{graph}.
\]
Mỗi thành phần chỉ xét các phần tử được $\OCLval$ chấp nhận; các liên kết đồ thị
là một tập ngữ nghĩa, trong đó các quan hệ trùng cho cùng liên kết UML được đồng
nhất. Cho $\liftNode$ ánh xạ từng điểm mọi thành phần có giá trị đối tượng qua
$\node$ và giữ nguyên các giá trị vô hướng. Khi đó
\[
(MM,M)\equiv_{val}G
\iff
\liftNode(\ObsObj(MM,M))=\ObsGraph(MM,G),
\]
với đẳng thức ngoại diên theo từng thành phần có kiểu. Bộ đọc bộ phận
$\Psi_{val}$ trả về chính xác bộ quan sát này và không tái dựng metadata UML
không phục vụ kiểm tra.
\end{definition}

Ngữ nghĩa nguồn được định nghĩa trên $\OCLval$:
\[
\sem{s}{OCL_{val}}(M,\rho):=
erase_S(\sem{s}{S}(M,tag_S(\rho))).
\]
Tính đầy đủ của parser cung cấp $decode_{val}(Parse(text(s)))=s$; phương trình
này nối biểu diễn text/AST với ngữ nghĩa nguồn mà không đồng nhất các siêu mô hình.

\section{Các phán đoán định kiểu tĩnh và binding}

Các kiểu tĩnh là
\[
\tau ::= Boolean\mid Integer\mid Real\mid String\mid C\mid Set(\tau).
\]
$\bot$ tại runtime không phải kiểu tĩnh. Quan hệ phù hợp $\leq$ gồm quan hệ phù
hợp lớp UML, $Integer\leq Real$, tính phản xạ và tính hiệp biến của tập hữu hạn.
Các bộ phân giải bộ phận tất định
\[
resolveClass,quad resolveAttribute,quad resolveNavigation,quad
resolveTypeOperation
\]
được giả thiết là đầy đủ đối với siêu mô hình và có tính hàm trên miền của
chúng; tra cứu mơ hồ hoặc không được hỗ trợ là không xác định.

Môi trường định kiểu $\Gamma$ là một stack và phép tra cứu chọn binding gần
nhất. Các phán đoán chính là
\[
MM;\Gamma\vdash b:\tau,
\qquad
MM;\Gamma\vdash ast\Rightarrow b:\tau.
\]
Kiểu kết quả số học và kiểu kết quả nhánh dùng các hàm bộ phận tất định. Cụ thể,
cộng, trừ và nhân số nguyên trả về $Integer$; chia số nguyên trả về $Real$; số
học trộn kiểu số trả về $Real$; và $join(\tau_1,\tau_2)$ là siêu kiểu chung nhỏ
nhất duy nhất khi tồn tại. Hàm kết quả không xác định làm admission thất bại.

Với literal tập và các phép hợp/giao tập, định nghĩa hàm bộ phận chuẩn
$setJoin$ trên kiểu phần tử như sau:
\[
\begin{array}{rcl}
setJoin(\tau,\tau)&=&\tau,\\
setJoin(Integer,Real)&=&setJoin(Real,Integer)=Real,\\
setJoin(C,D)&=&\operatorname{lub}_{MM}(C,D)
\quad\text{khi siêu lớp chung nhỏ nhất là duy nhất}.
\end{array}
\]
Hàm nhiều đối số, chỉ dùng cho danh sách không rỗng, là phép gấp trái của các
phương trình trên và chỉ xác định khi mọi bước trung gian xác định. Không có
phương trình giữa kiểu nguyên thủy không tương thích, giữa kiểu nguyên thủy và
lớp, hoặc cho kiểu collection lồng nhau.

Khi $\tau\leq\upsilon$ là một quan hệ nâng kiểu được $setJoin$ tạo ra, định
nghĩa phép nhúng chuẩn có kiểu
\[
\iota^I_{\tau\hookrightarrow\upsilon}:
Val_I^\bot(\tau)\longrightarrow Val_I^\bot(\upsilon)
\]
như sau: nó là đồng nhất khi $\tau=\upsilon$; ánh xạ Integer sang giá trị
Real64 chính xác khi $\tau=Integer,\upsilon=Real$; giữ nguyên định danh thực
thể hoặc nút khi $\tau=C,\upsilon=D$ và $C$ phù hợp với $D$; và ánh xạ
$\bot$ sang $\bot$. Điều kiện $ScalarClosed$ bảo đảm phép nâng Integer--Real
được xác định. Với $S:Set(\tau)$, đặt
\[
\operatorname{upSet}^{I}_{\tau\hookrightarrow\upsilon}(S)
=\{\iota^I_{\tau\hookrightarrow\upsilon}(v)\mid
v\in\finiteSet(S)\}.
\]
Ngữ nghĩa Set literal, $union$ và $intersection$ trước hết nâng mọi phần tử
vào kiểu kết quả $\upsilon=setJoin(\ldots)$; giao tập và loại trùng sau đó dùng
$\eqval^{\upsilon}$. Vì vậy không có so sánh không kiểu giữa hai miền phần tử.
Các phép $setJoin$, $\iota$ và $\operatorname{upSet}$ là helper định kiểu của
các hàm ngữ nghĩa SF-01--SF-08, không định nghĩa thêm một bộ đánh giá độc lập.

\begin{lemma}[M3a Tính hàm của $setJoin$]
Nếu $setJoin(\tau_1,\ldots,\tau_n)=\sigma_1$ và
$setJoin(\tau_1,\ldots,\tau_n)=\sigma_2$, thì $\sigma_1=\sigma_2$.
Khi được xác định, kết quả là siêu kiểu chung nhỏ nhất chuẩn duy nhất của các
kiểu phần tử. Nếu kết quả là $\upsilon$, thì mỗi
$\iota^I_{\tau_i\hookrightarrow\upsilon}$ là một hàm toàn phần trên các giá
trị được chấp nhận và $\operatorname{upSet}^{I}$ cũng có tính hàm.
\end{lemma}
\begin{proof}
Mỗi phương trình cơ sở có các miền rời nhau và trả về một kết quả duy nhất.
Trường hợp lớp dùng $\operatorname{lub}_{MM}$ chỉ khi phần tử nhỏ nhất là duy
nhất; nếu có hai ứng viên tối tiểu không so sánh được thì hàm không xác định.
Quy nạp trên số đối số của phép gấp trái, dùng tính hàm ở bước trước và của
phương trình nhị phân, cho $\sigma_1=\sigma_2$. Tính chất siêu kiểu và tính nhỏ
nhất được bảo toàn qua từng bước gấp.
Với phép nhúng, phân trường hợp theo phương trình cuối cùng của $setJoin$.
Trường hợp đồng nhất
là trực tiếp. Trường hợp Integer--Real dùng tính chính xác của phép đổi sang
binary64 do $ScalarClosed$ yêu cầu, nên hai phía nhận cùng giá trị Real64 hữu
hạn. Trường hợp lớp chỉ thay chú thích kiểu tĩnh và giữ nguyên thực thể hoặc
nút; bottom được giữ nguyên. Đây là toàn bộ miền của $setJoin$, nên phép nhúng
có tính hàm và toàn phần trên miền được chấp nhận. Tính hàm của
$\operatorname{upSet}^{I}$ suy ra bằng phép dựng ảnh hữu hạn.
\end{proof}

Phán đoán thứ hai được định hướng theo cú pháp và định nghĩa
$T_{\mathrm{BIND}}$. $resolveProperty$ là bộ phân loại rời nhau có tính hàm,
trả về một thuộc tính hoặc một bộ navigation. Bảng sau là đầy đủ đối với
$\OCLval$; $e\Rightarrow b:\tau$ viết tắt cho
$MM;\Gamma\vdash e\Rightarrow b:\tau$.

\begingroup\small
\begin{longtable}{p{0.12\textwidth}p{0.45\textwidth}p{0.35\textwidth}}
\toprule
Luật & Tiền đề & Kết quả/kiểu Bound\\
\midrule
B-Inv & $resolveClass(ctx)=C$; bind $p$ thành Boolean dưới $\{self:C\}$ &
$ContextInvariant(ctx,R,p)\Rightarrow BInvariant(C,R,b)$\\
B-Self & $\Gamma(self)=C$ & $self\Rightarrow BSelf(C):C$\\
B-Var & $nearest(\Gamma,x)=\tau$ & $x\Rightarrow BVar(x,\tau):\tau$\\
B-Lit & $typeOfLiteral(c)=\tau$ & $c\Rightarrow BLiteral(c,\tau):\tau$\\
B-SetLit & $n\geq1$; $e_i\Rightarrow b_i:\tau_i$; $setJoin(\tau_1,\ldots,\tau_n)=\tau$;
không $\tau_i$ nào là collection & $Set\{e_1,\ldots,e_n\}\Rightarrow BSetLiteral([b_1,\ldots,b_n]):Set(\tau)$\\
B-Attr & $e\Rightarrow b:D$; $resolveProperty(D,n,[])=Attr(a,\tau)$ &
$e.n\Rightarrow BAttribute(b,a,\tau):\tau$\\
B-Nav & $e\Rightarrow b:\sigma$, $\sigma\in\{D,Set(D)\}$;
$q_i\Rightarrow b_i:\sigma_i$; resolver trả về $(A,fr,tr,dir,E,\bar\sigma)$ &
$e.r[\bar q]\Rightarrow BNavigation(b,A,fr,tr,dir,\bar b):Set(E)$\\
B-AllInst & $resolveClass(name)=C$ & $C.allInstances()\Rightarrow BAllInstances(C):Set(C)$\\
B-Not & $p\Rightarrow b:Boolean$ & $not\ p\Rightarrow BNot(b):Boolean$\\
B-Bool & $e_i\Rightarrow b_i:Boolean$; $op\in\{and,or,xor,implies\}$ &
$e_1\ op\ e_2\Rightarrow BBinary(op,b_1,b_2):Boolean$\\
B-Eq & $e_i\Rightarrow b_i:\tau_i$; các kiểu so sánh được; $op\in\{=,<>\}$ &
$BBinary(op,b_1,b_2):Boolean$\\
B-Ord & $\tau_i$ tương thích số; $op\in\{<,\leq,>,\geq\}$ &
$BBinary(op,b_1,b_2):Boolean$\\
B-Arith & $arithResult(op,\tau_1,\tau_2)=\tau$ &
$BBinary(op,b_1,b_2):\tau$\\
B-If & $c:Boolean$, các nhánh $\tau_1,\tau_2$; $join(\tau_1,\tau_2)=\tau$ &
$BIf(b_c,b_t,b_f,\tau):\tau$\\
B-Let & $init\Rightarrow b_i:\tau_1$, $\tau_1\leq T$; thân $\tau_2$ dưới $\Gamma[x:T]$ &
$BLet(x,T,b_i,b_b,\tau_2):\tau_2$\\
B-Exists/ForAll & $S:Set(\tau)$; $P:Boolean$ dưới $\Gamma[x:\tau]$ &
$BIterator(op,b_S,x,b_P):Boolean$\\
B-Select/Reject & cùng tiền đề nguồn và thân Boolean &
$BIterator(op,b_S,x,b_P):Set(\tau)$\\
B-Collect & $S:Set(\tau)$; $E:\sigma$ dưới $\Gamma[x:\tau]$; $\sigma\neq Set(\_)$ &
$BIterator(collect,b_S,x,b_E):Set(\sigma)$\\
B-IsUnique & $S:Set(\tau)$; $E:\sigma$ dưới $\Gamma[x:\tau]$; $\sigma\neq Set(\_)$ &
$BIterator(isUnique,b_S,x,b_E):Boolean$\\
B-Member & $S:Set(\tau)$; $e:\sigma$; các kiểu so sánh được &
$BCollectionOp(includes/excludes,b_S,[b_e]):Boolean$\\
B-SetRel & $S:Set(\tau)$; $T:Set(\sigma)$; các kiểu phần tử so sánh được &
$BCollectionOp(includesAll/excludesAll,b_S,[b_T]):Boolean$\\
B-SetComb & $S:Set(\tau)$; $T:Set(\sigma)$; $setJoin(\tau,\sigma)=\upsilon$ &
$BCollectionOp(union/intersection,b_S,[b_T]):Set(\upsilon)$\\
B-AsSet & $S:Set(\tau)$ & $BCollectionOp(asSet,b_S,[]):Set(\tau)$\\
B-Size & $S:Set(\tau)$ & $BCollectionOp(size,b_S,[]):Integer$\\
B-Empty & $S:Set(\tau)$ & $BCollectionOp(isEmpty/notEmpty,b_S,[]):Boolean$\\
B-Kind & $e:D$; $resolveClass(name)=C$ & $BTypeOp(kindOf,b,C):Boolean$\\
B-Cast & $e:D$; $resolveClass(name)=C$ & $BTypeOp(cast,b,C):C$\\
\bottomrule
\end{longtable}
\endgroup

Mọi phán đoán con bị lược bỏ trong một dòng rút gọn đều là tiền đề dưới
$\Gamma$ hiện tại. Chỉ thân iterator và let dùng mở rộng được hiển thị. Nếu bất
kỳ tiền đề nào thất bại thì binder bộ phận không xác định. $TopLevelQuery$
không phải gốc invariant và nằm ngoài Định lý~6. \texttt{oclIsTypeOf} là phần
mở rộng của prototype/parser và không có luật binder được chứng nhận.

\subsection{Chứng nhận ngữ nghĩa Source/Bound đầy đủ}

Đánh giá Source và Bound được định nghĩa trên các miền có gắn thẻ rời nhau
$Val_S$ và $Val_B$. Các họ phép toán của chúng
$truth_S,finite_S,eq_S,ord_S,arith_S,nav_S$ và
$truth_B,finite_B,eq_B,ord_B,arith_B,nav_B$ được định nghĩa độc lập từ cú pháp
và thông tin định kiểu tương ứng; không bộ đánh giá nào gọi ngữ nghĩa VA. Viết
$S(a,\rho_S)=\sem{a}{S}(M,\rho_S)$ và
$B(b,\rho_B)=\sem{b}{B}(M,\rho_B)$.

Định nghĩa quan hệ logic có kiểu nhỏ nhất $\sim_{SB}$ bằng các thẻ Boolean, vô
hướng, thực thể, bottom và tập hữu hạn tương ứng, trong đó các tập hữu hạn liên
hệ theo ngoại diên. Các môi trường liên hệ theo từng điểm. Định nghĩa các phép
xóa thẻ toàn phần
\[
erase_S:Val_S\to Val_{obj},\qquad erase_B:Val_B\to Val_{obj}
\]
theo từng thành phần trên bottom, các thẻ Boolean/vô hướng/thực thể và đệ quy
trên tập hữu hạn. Cầu nối xóa thẻ là
\[
v_S\sim_{SB}v_B\Longrightarrow erase_S(v_S)=erase_B(v_B).
\]
Với môi trường không gắn thẻ $\rho$, cho $tag_S(\rho)$ và $tag_B(\rho)$ áp dụng
các thẻ tương ứng theo từng điểm, rồi định nghĩa ngữ nghĩa công khai bởi
\[
\sem{s}{OCL_{val}}(M,\rho)
=erase_S(\sem{s}{S}(M,tag_S(\rho))),
\]
\[
\sem{b}{Bound}(M,\rho)
=erase_B(\sem{b}{B}(M,tag_B(\rho))).
\]
Bảng sau bao phủ mọi dòng binder thành công; ngữ nghĩa con tương ứng sinh ra
ngữ nghĩa cha tương ứng.

\begingroup\scriptsize
\begin{longtable}{p{0.20\textwidth}p{0.35\textwidth}p{0.38\textwidth}}
\toprule Dạng Source/Bound & Phương trình Source & Phương trình Bound / lý do\\ \midrule
self, biến, literal & giá trị môi trường hoặc literal gắn thẻ Source & giá trị gắn thẻ Bound tương ứng\\
Set literal & nâng từng phần tử Source tới kiểu $setJoin$ rồi lấy tập ngoại diên & cùng phép nâng có kiểu trên các phần tử Bound tương ứng\\
thuộc tính & quan sát thuộc tính Source trên $S(src,\rho_S)$ & quan sát thuộc tính Bound trên $B(b,\rho_B)$\\
navigation có qualifier & hợp của $nav_S$ trên các receiver gắn thẻ Source sau khi đánh giá qualifier có thứ tự & hợp của $nav_B$ với liên kết, role, hướng, số ngôi và thứ tự qualifier đã phân giải\\
$allInstances(C)$ & $allInstances_{obj}(C)$ & cùng tập lớp đã phân giải\\
not/Boolean nhị phân & $truth_S/bool_S$ được định nghĩa độc lập & $truth_B/bool_B$ được định nghĩa độc lập\\
đẳng thức/thứ tự & $eq_S/ord_S$ trên các thẻ Source & $eq_B/ord_B$ dùng các kiểu tĩnh đã phân giải\\
số học & $arith_S(op,S(e_1,\rho_S),S(e_2,\rho_S))$ & $arith_B(op,B(b_1,\rho_B),B(b_2,\rho_B))$\\
if & nhánh được $truth_S$ chọn & nhánh tương ứng được $truth_B$ chọn\\
let & mở rộng môi trường Source & mở rộng môi trường Bound tương ứng\\
exists/forAll & lượng hóa trên $finite_S$ & lượng hóa trên $finite_B$ dưới các mở rộng tương ứng\\
isUnique & tính đơn ánh của ảnh thân trên nguồn hữu hạn & cùng phép kiểm tra đơn ánh hữu hạn trên thân Bound\\
select/reject & giữ phần tử có thân đúng/không đúng theo kiểm tra & cùng phép bao hàm tập hữu hạn cho thân Bound\\
collect theo tập hữu hạn & $\{S(E,\rho[x\mapsto s])\mid s\in\finiteSet(S(src,\rho))\}$ & cùng ảnh dùng $B(b_E,\rho[x\mapsto s])$; không làm phẳng hoặc giữ bội số\\
quan hệ thuộc/quan hệ tập & quan hệ thuộc, tập con hoặc rời nhau trên $\finiteSet$ & cùng quan hệ tập hữu hạn có kiểu\\
union/intersection/asSet & nâng về kiểu kết quả rồi hợp, giao hoặc đồng nhất ngoại diên & cùng phép nâng và phép toán tập hữu hạn trên Bound\\
size/tính rỗng & lực lượng hoặc tính rỗng của $\finiteSet(S(src,\rho))$ & cùng phép toán trên $\finiteSet(B(b_S,\rho))$\\
kindOf/cast & quan hệ phù hợp hoặc giá trị/bottom có guard cho $C$ đã phân giải & phép toán Bound tương ứng được định nghĩa độc lập\\
gốc invariant & $\boolval(S(p,\rho[self\mapsto o]))$ cho từng đối tượng ngữ cảnh & $\boolval(B(b,\rho[self\mapsto o]))$\\
\bottomrule
\end{longtable}
\endgroup

\begin{lemma}[SB0 Tương ứng toán tử độc lập]
Với mọi dòng định kiểu được chấp nhận, các toán hạng tương ứng sinh kết quả
tương ứng dưới các phép toán Source và Bound được định nghĩa độc lập. Mở rộng
môi trường bằng các giá trị tương ứng bảo toàn quan hệ theo từng điểm.
\end{lemma}
\begin{proof}
Phân tích hữu hạn các trường hợp trên thẻ Boolean/vô hướng được chấp nhận và
lập luận ngoại diên cho tập hữu hạn. Set literal, hợp và giao dùng cùng phép
nhúng chuẩn có tính hàm từ M3a trước khi so sánh ngoại diên. Navigation
dùng metadata đã phân giải duy nhất của B3; mở rộng iterator và let dùng B4.
Không constructor VA nào xuất hiện trong chứng minh này.
\end{proof}
Do đó bảng là phân hoạch trường hợp cục bộ đầy đủ cho Định lý~1.

\section{Danh mục bổ đề}

\subsection{Các bổ đề biểu diễn cho Định lý 0}

\begin{lemma}[R1 Tính đơn ánh của đối tượng]
Với mọi $o_1,o_2\in \Obj(M)$,
\[
\node(o_1)=\node(o_2) \iff o_1=o_2.
\]
\end{lemma}
\begin{proof}
Theo hợp đồng mã hóa, mỗi đối tượng được gán một định danh ổn định, đơn ánh và
chính xác một nút đồ thị lưu định danh đó. Nếu hai nút đối tượng bằng nhau thì
chúng có cùng \texttt{use\_id}; tính đơn ánh của định danh đối tượng cho
$o_1=o_2$. Ngược lại, nếu $o_1=o_2$ thì tính hàm của phép dựng nút cho
$\node(o_1)=\node(o_2)$ sau khi thay hai đối số bằng cùng một đối tượng.
\end{proof}

\begin{lemma}[R2 Tính chính xác đối tượng--nút]
Với mọi $o\in \Obj(M)$, tồn tại duy nhất $n\in \Node(G)$ sao cho
\[
\RepG(o,n).
\]
Ngược lại, với mọi nút đối tượng nhìn thấy bởi kiểm tra
$n\in\ObjNodeVal(G)$, tồn tại duy nhất $o\in\Obj(M)$ sao cho
$\RepG(o,n)$. Tương đương,
\[
\ObjNodeVal(G)=\{\node(o)\mid o\in\Obj(M)\}.
\]
Ký hiệu $\node(o)$ chỉ nút duy nhất này; do đó
$\prop(\node(o),\texttt{use\_id})=\id(o)$.
\end{lemma}
\begin{proof}
Theo mệnh đề kiến tạo đối tượng--nút, $\PhiEnc$ sinh ít nhất một nút từ mỗi đối
tượng và lưu $\id(o)$ làm \texttt{use\_id}. Mệnh đề mỗi đối tượng một nút cho
tính duy nhất. Mệnh đề không có nút đối tượng giả cung cấp sự tồn tại và duy
nhất theo chiều ngược cho mọi nút trong $\ObjNodeVal(G)$. Vì vậy $\node$ là
song ánh giữa $\Obj(M)$ và $\ObjNodeVal(G)$; đẳng thức accessor suy ra từ định
nghĩa $\RepG$.
\end{proof}

\begin{lemma}[R3 Bảo toàn kiểu]
Với mọi đối tượng $o$ và lớp $C$,
\[
o \ \mathrm{conformsTo}\ C
\iff
\node(o)\in I^{graph}_{class}(C).
\]
Tương đương,
\[
o \ \mathrm{conformsTo}\ C
\iff
(\node(o))-[:\mathrm{ObjectInstanceOf}]\to(\classNode(C)).
\]
\end{lemma}
\begin{proof}
Chiều thuận đúng vì phép mã hóa vật chất hóa các cạnh
\texttt{ObjectInstanceOf} cho lớp runtime và mọi siêu kiểu phù hợp. Chiều
ngược suy ra từ điều kiện không có cạnh kiểu giả: mọi cạnh kiểu nhìn thấy bởi
kiểm tra đều được sinh từ một sự kiện phù hợp thực sự.
\end{proof}

\begin{lemma}[R4 Bảo toàn thuộc tính]
Với mọi đối tượng $o$ và thuộc tính $a$,
\[
\mathrm{attrVal}_M(o,a)=I^{graph}_{attr}(\node(o),a).
\]
\end{lemma}
\begin{proof}
Đây chính là mệnh đề accessor thuộc tính của hợp đồng mã hóa đồ thị: accessor
đồ thị chuẩn $I^{graph}_{attr}$ quan sát cùng giá trị với phép tra cứu thuộc
tính trên mô hình đối tượng.
\end{proof}

\begin{lemma}[R5 Bảo toàn liên kết]
Với mọi liên kết nhị phân $A$,
\[
(o_1,sr,\bar q_s,o_2,tr,\bar q_t)\in \mathrm{links}_M(A)
\]
khi và chỉ khi tồn tại $r\in \Rel(G)$ sao cho:
\[
\src(r)=\node(o_1),\quad \trg(r)=\node(o_2),
\]
\[
\prop(r,\texttt{associationKey})=associationKey_{MM}(A),\quad
\prop(r,\texttt{sourceRole})=sr,\quad
\prop(r,\texttt{targetRole})=tr,
\]
và hai danh sách qualifier đã mã hóa khớp theo đúng đầu:
\[
\mathrm{qualifierMatches}(\mathrm{encodeQualifierList}(\bar q_s),
  \prop(r,\texttt{sourceQualifiers}))
\]
\[
\land\quad
\mathrm{qualifierMatches}(\mathrm{encodeQualifierList}(\bar q_t),
  \prop(r,\texttt{targetQualifiers})).
\]
\end{lemma}
\begin{proof}
Nếu liên kết UML tồn tại, $\PhiEnc$ xuất ít nhất một quan hệ \texttt{Link*}
tương ứng có cùng \texttt{associationKey}, role, đầu mút và hai danh sách
qualifier có hướng. Ngược lại,
theo điều kiện không có liên kết giả, mọi quan hệ \texttt{Link*} nhìn thấy bởi
kiểm tra đều được sinh từ một liên kết UML. Tính duy nhất của đầu mút suy ra từ
R1 và R2. Chiều phản xạ của đẳng thức qualifier dùng tính đơn ánh của
$\mathrm{encodeQualifierList}$ trên các danh sách có thứ tự, đúng số ngôi đã
khai báo và thuộc các loại qualifier nguyên thủy được hỗ trợ. Không cần tính
duy nhất của quan hệ vì các hiện thực hóa phía sau dùng ngữ nghĩa tập.
\end{proof}

\begin{lemma}[R6 Bảo toàn navigation]
Với mọi navigation nhị phân được hỗ trợ có các biểu thức qualifier
$\bar q=[q_1,\ldots,q_k]$, đặt
$rawQs=[\sem{q_1}{obj}(\rho),\ldots,\sem{q_k}{obj}(\rho)]$ và
$encodedQs=\mathrm{encodeQualifierList}(rawQs)$. Khi đó
\[
\{\node(y)\mid y\in nav_{obj}(o,A,fromRole,toRole,direction,rawQs)\}
=
nav_{graph}(\node(o),A,fromRole,toRole,direction,encodedQs).
\]
Theo quy ước association-end của OCL/USE, navigation thuận đối chiếu $rawQs$
với qualifier của đầu nguồn và trường \texttt{sourceQualifiers}; navigation
ngược đối chiếu qualifier của đầu đích và trường \texttt{targetQualifiers}.
Quy ước này là một phần của metadata hướng đã phân giải, không được suy ra từ
tên role tại thời gian render.
\end{lemma}
\begin{proof}
Hai bao hàm suy ra từ R5. Một kết quả navigation UML tương ứng với liên kết
được xuất thành quan hệ đồ thị, nên nút đích của nó xuất hiện trong
$nav_{graph}$. Ngược lại, mọi kết quả navigation đồ thị đến từ một quan hệ
\texttt{Link*} khớp; theo điều kiện không có liên kết giả, quan hệ đó tương ứng
với một liên kết UML. Metadata navigation đã phân giải cố định việc khớp hướng
và role. Các biểu thức qualifier được đánh giá trước; sau đó phép so sánh dùng
phép tuần tự hóa danh sách đơn ánh, theo từng thành phần và bảo toàn thứ tự.
\end{proof}

\begin{lemma}[R7 Bảo toàn allInstances]
Với mọi lớp $C$,
\[
\{\node(o)\mid o\in allInstances_{obj}(C)\}
=
allInstances_{graph}(C).
\]
\end{lemma}
\begin{proof}
Theo ngữ nghĩa đối tượng hữu hạn, $o\in allInstances_{obj}(C)$ khi và chỉ khi
$o$ phù hợp với $C$. Theo R3, điều này tương đương với việc $\node(o)$ có cạnh
\texttt{ObjectInstanceOf} tới $\classNode(C)$, tức chính xác là quan hệ thuộc
$allInstances_{graph}(C)$.
\end{proof}

\begin{remark}[Vị thế kiến tạo của $\Phi$]
Định lý~0 là tính đầy đủ dưới hợp đồng mã hóa tường minh, không phải định lý về
mọi bộ mã hóa tùy ý. Định nghĩa $\Phi^{\star}(MM,M)$ là đồ thị nhỏ nhất nhìn
thấy bởi kiểm tra, được sinh bởi: một schema subgraph $G_{M_2}$ lưu đầy đủ các
khai báo M2 cần cho resolver được chấp nhận, gồm lớp, thuộc tính và kiểu, liên
kết nhị phân và đầu mút, role, qualifier cùng kế thừa; chính xác một nút
\texttt{Class} có
$classKey$ đơn ánh cho mỗi lớp; chính xác một nút đối tượng có
\texttt{use\_id} đơn ánh cho mỗi đối tượng; chính xác một cạnh
\texttt{ObjectInstanceOf} cho mỗi cặp phù hợp đúng; chính xác một nút
\texttt{AttributeValue} và cạnh \texttt{ObjectHasAttribute} cho mỗi vị trí
$(o,a)$ được chấp nhận, dùng khóa $attributeKey_{MM}(a)$ và lưu
$storeScalar(attrVal_M(o,a))$; cùng chính xác một quan hệ \texttt{Link} cụ thể
cho mỗi liên kết UML nhị phân được chấp nhận, lưu $associationKey_{MM}(A)$,
các role và hai danh sách \texttt{sourceQualifiers}/\texttt{targetQualifiers}
có thứ tự. Không có sự kiện nào khác nhìn thấy
bởi kiểm tra. Ở đây $storeScalar$ đơn ánh và
$readScalar(storeScalar(v))=v$. Các chiều thuận của R2--R5 suy ra từ phép sinh,
các chiều ngược suy ra từ tính nhỏ nhất/chính xác; R6--R7 suy ra theo ngoại diên.
\end{remark}

\begin{definition}[Tính đầy đủ của adapter mã hóa cụ thể]
Với đồ thị cụ thể $G_{impl}$, profile adapter $P$ cung cấp các quan sát về đối
tượng/định danh, lớp, kiểu, thuộc tính, navigation, all-instances và một
$GraphMMView_P$ trên schema subgraph M2.
$AdapterAdequate(P,G_{impl},MM,M)$ đúng khi và chỉ khi mọi quan sát có thể đạt
tới bởi định lý đều được định nghĩa và bằng quan sát tương ứng của
$\Phi^{\star}(MM,M)$ sau ánh xạ đối tượng--nút bảo toàn định danh. Vì vậy,
prototype hiện dùng tra cứu chính xác \texttt{classKey},
\texttt{attributeKey}, \texttt{associationKey}, cùng
\texttt{sourceQualifiers}/\texttt{targetQualifiers} theo hướng; các quan sát đó
phải được chứng minh đầy đủ, không giả và duy nhất trong profile đã công bố.
Vì production sinh text trực tiếp, bridge implementation phải kiểm tra text đó
bằng parser độc lập, expected formal tree và parser Neo4j đã ghim; nó không được
giả định production đi qua \texttt{RawCypherAst}. Cho tới khi các nghĩa vụ phổ
quát còn lại được thỏa, Định lý~0--6 áp dụng cho profile chuẩn, không tự động áp
dụng cho toàn bộ prototype.
\end{definition}

\begin{remark}[Ranh giới binding và thực thi]
Binder hoặc adapter siêu mô hình được phép duyệt các relationship trong
$G_{M_2}$ để phân giải association end, role, hướng ngữ nghĩa, kiểu đích,
qualifier và kế thừa. Nó không được duyệt các relationship link của $G_{M_1}$
để đoán ý nghĩa property hoặc chuyên biệt hóa truy vấn theo dataset hiện tại.
Kết quả binding chứa metadata M2 đã phân giải; $T_{\mathrm{TEXT}}$ dùng metadata
đó cùng profile mã hóa để sinh pattern, và chỉ khi thực thi pattern mới khớp
các relationship instance của M1.
\end{remark}

\begin{remark}[Trạng thái adapter của prototype]
Prototype hiện tại cung cấp bằng chứng bộ phận, chưa hoàn thành nghĩa vụ
$AdapterAdequate$. Adapter chuẩn dùng khóa định tính cho lớp, thuộc tính và liên
kết, lưu role cùng qualifier có thứ tự theo cả hai hướng, và có thể dựng M2 view
từ siêu mô hình USE hoặc từ phần M2 trong graph. Oracle snapshot độc lập không
tìm thấy fact R1--R7 thiếu hoặc giả trên ba fixture khác nhau; chín mutation về
định danh, kiểu, thuộc tính, link, hướng qualifier, all-instances và khóa đều bị
phát hiện. Trên fixture kế thừa có qualifier, hai M2 view bằng nhau ngoại diên
và sinh cùng Cypher cùng ánh xạ tham số. Tuy nhiên, đây vẫn là bằng chứng hữu hạn
đa fixture, không phải chứng minh phổ quát của $AdapterAdequate$.

Production sinh text trực tiếp từ query plan. Parser độc lập đã chuẩn hóa ổn
định 47 truy vấn admitted; witness bao phủ 17/17 constructor plan; 52 truy vấn
đã khớp chính xác canonical token/group tree được duyệt, gồm nested shadowing.
Parser Neo4j Cypher 5 từ artifact \texttt{cypher-parser-factory-2026.06.0.jar}
cũng chấp nhận 52/52 và selected internal-AST-kind projection khớp các witness;
direction, \texttt{COUNT}/\texttt{COLLECT} và nested scope/alias được giữ. Đây
là bằng chứng hữu hạn, chưa chứng minh full-AST, semantic-analysis hoặc tương
đương phổ quát với raw AST hình thức.

Từ ngày 2026-08-01, document và single-rule validation dùng
\texttt{compileFileWithCertifiedContextInvariants}; mọi \texttt{Class::inv}
được compile đều đi qua pre-admission, certified binder và post-binding
admission như instrumented path. Trước execution, service đối chiếu lại
text/parameter với instrumented result và chạy expression-level
\texttt{ScalarClosed} bằng exact \texttt{attributeKey} observations từ graph.
Research Tool chạy graph \texttt{BottomSeparated}, stored-scalar scan,
USE-observed expression \texttt{ScalarClosed} và generated-bottom gate. Full
static gate sau integration đạt 247/247, mutation score 20/20. Runtime
follow-up trên Neo4j Kernel 2026.06.0 Enterprise/Cypher 5/database
\texttt{demo} đạt 47/47 violation sets với premise gate mới; năm dialect smoke probe
và ba bottom/scalar/receiver probe đều pass, không skip. Kết quả này không thay
thế ngữ nghĩa hình thức đầy đủ của Neo4j. Chứng minh giấy hiện cũng chưa được cơ
giới hóa trong proof assistant.
\end{remark}

\begin{lemma}[PA14d: Đồng thuận projection với parser Neo4j đã chọn]
Gọi $parseN5$ là \texttt{AstParserFactory(CypherVersion.Cypher5)} phiên bản
2026.06.0 và $projN5$ là phép chiếu AST nội bộ bỏ source position/object identity
nhưng giữ constructor/field order, schema/property/function name, parameter,
literal, direction và alias. Với 52 trường hợp đã kiểm tra (gồm 47 certified),
$parseN5(text(P_i))$ được định nghĩa. Mọi yêu cầu AST-kind của witness theo
17/17 constructor đều có đồng thời một quan sát từ parser fragment độc lập và
một node/type tương ứng trong $projN5(parseN5(text(P_i)))$.
\end{lemma}
\begin{proof}
\texttt{Neo4jCypherAstBridgeTest} ghim đúng JAR, parse 52/52, kiểm projection lặp
lại là deterministic, từ chối hai syntax mutant và phân biệt hai direction cùng
\texttt{CountExpression}/\texttt{CollectExpression}. Fixture shadowing có hai
\texttt{FilterScope} và giữ alias tường minh. \texttt{ExpectedFormalCypherTreeVerifier}
chỉ chấp nhận AST-kind khi cả hai oracle cùng witness. Focused suite đạt 6/6;
checkpoint static khi đó đạt 250/250 cả bình thường và Maven offline, không
failure/error/skip. Đây là chứng minh kiểm
tra hữu hạn của selected projection, không phải full-AST hoặc ngữ nghĩa Neo4j.
\end{proof}

\begin{remark}[PO-15: bằng chứng runtime một-hàng-mỗi-CY]
Ngày 2026-08-01, \texttt{Cypher5ValDialectRealNeo4jTest} chạy đúng một query,
fixture và oracle cho mỗi CY1--CY9. Fixture transaction có hai physical link
paths tới một semantic target, hai hướng navigation, exact
\texttt{associationKey}/role/qualifier, bottom được biểu diễn, null, duplicate,
alias tường minh và nhánh tương quan lồng nhau. Kết quả 9/9 pass trên Neo4j
Kernel 2026.06.0 Enterprise, Cypher 5, database \texttt{demo}, driver dependency
5.21.0. Manifest TSV ghim profile và SHA-256 của renderer, canonical encoding,
probe matrix và runtime harness; guard 2/2 bắt duplicate/FAIL row, altered
observation và stale hash.
Checkpoint static sau đó đạt 252/252 từ Maven cache, mutation score
20/20, proof sync pass. Manifest đó hiện là schema-v1/contract-.2, dirty và bị
freshness guard của contract-.3 từ chối; vì vậy PO-15 đã mở lại cho clean capture.
CY1--CY9 vẫn là tiên đề hành vi tin cậy chứ không thành bổ đề toán học.
\end{remark}

\begin{remark}[PO-16: non-vacuity và exact-ID differential]
Manifest kỳ vọng được duyệt trước khi chạy Cypher cố định exact violation IDs
và lý do ngữ nghĩa cho C01--C47. Fixture hợp lệ multiplicity tạo đồng thời
satisfying và violating objects cho cả 28 case có countermodel. Native USE và
Neo4j Kernel 2026.06.0 cùng trả đúng các tập đã chốt: 47/47 equivalent, 28 mixed,
19 all-pass, không all-violate hoặc empty context. Mười chín all-pass là các
profile tautology: identity/reflexivity (C01,C02,C16,C24), literal/order/control
(C03,C15,C26), multiplicity/context/type (C09,C11,C33,C34,C40,C41), arithmetic
không overflow dưới \texttt{ScalarClosed} (C22,C23), và luật cardinality/tập hữu
hạn (C30,C31,C35,C37). Guard evidence ghim expected manifest, fixture, renderer,
encoding, runtime harness và profile; checkpoint static khi đó đạt 255/255,
mutation score 20/20. Manifest hiện là schema-v1/contract-.2, dirty và stale,
nên PO-16 mở cho clean schema-v2 recapture. Đây vẫn là evidence hữu hạn, không phải equivalence cho
mọi model/OCL input.
\end{remark}

\subsection{Các bổ đề biến đổi và hiện thực hóa}

\begin{lemma}[F1 Biểu diễn $\OCLval$ và rebinding Bound]
Với mọi $s\in\OCLval$, $encode_{ast}(s)\in\ASTval$ và
$decode_{val}(encode_{ast}(s))=s$. Nếu binding thành công với
$T_{\mathrm{BIND}}(encode_{ast}(s),MM)=v$, thì:
\begin{enumerate}[label=(\roman*)]
  \item $encode_{ast}(s)$ phù hợp siêu mô hình OCL AST;
  \item $v\in\BoundOCLval(MM)$;
  \item $erase_{bound}(v)\equiv_{\alpha}encode_{ast}(s)$;
  \item mọi construct trong $s$ và $v$ thuộc phạm vi định lý;
  \item không construct bị loại nào xuất hiện.
\end{enumerate}
\end{lemma}
\begin{proof}
M1 chứng minh mệnh đề Source--AST. Với round trip Bound, quy nạp trên suy dẫn
binding thành công. Constructor cơ sở xóa thành nút parser tương ứng; constructor
hợp thành dùng giả thiết quy nạp cho các phần tử con. B1 khôi phục tham chiếu UML,
B3 khôi phục tuple navigation duy nhất và B4 khôi phục phạm vi iterator/let.
Hai cây chỉ có thể khác bởi tên mới, nên tương đương alpha. Construct bị loại
không có phương trình decode hoặc trường hợp binding.
\end{proof}

\begin{lemma}[B1 Tính đúng đắn của phân giải tên]
Nếu $MM;\Gamma\vdash ast\Rightarrow b:\tau$, mọi tham chiếu biến biểu thị
binding gần nhất trong $\Gamma$, và mọi tham chiếu UML được lưu là phần tử duy
nhất do resolver đầy đủ trả về. Vì vậy nó biểu thị cùng thực thể ngữ nghĩa với
AST nguồn được chấp nhận.
\end{lemma}
\begin{proof}
Quy nạp trên suy dẫn binding. Các trường hợp biến và self dùng tra cứu gần nhất
trên stack. Các trường hợp lớp, thuộc tính, navigation và toán tử kiểu dùng
tính đầy đủ của resolver và lưu phần tử UML được trả về thay vì tên bề mặt.
Các trường hợp phức hợp dùng giả thiết quy nạp; tra cứu không được hỗ trợ không
có suy dẫn.
\end{proof}

\begin{lemma}[B2 Tính đúng đắn của gán kiểu]
Nếu $MM;\Gamma\vdash ast\Rightarrow b:\tau$ thì
$MM;\Gamma\vdash b:\tau$. Mọi suy dẫn con có kiểu đều có tính chất tương tự;
khi đánh giá của nó cho giá trị khác bottom, giá trị đó thuộc diễn giải của
kiểu đã gán.
\end{lemma}
\begin{proof}
Quy nạp trên suy dẫn binding. Mọi luật binding kết luận bằng luật định kiểu
tương ứng đã nêu. Một quy nạp đồng thời chứng minh quan hệ thuộc của giá trị
bằng tính đầy đủ kiểu của resolver, tính đóng của phép toán vô hướng, phép dựng
tập hữu hạn và các giả thiết quy nạp dưới mở rộng môi trường.
\end{proof}

\begin{lemma}[B3 Tính tất định của phân giải navigation]
Nếu binding một lời gọi thuộc tính được hỗ trợ thành navigation thành công, nó
sinh một bộ duy nhất
\[
(A,fromRole,toRole,direction,qualifierTypes)
\]
tương thích với kiểu receiver và $MM$. Navigation mơ hồ hoặc không được hỗ trợ
bị từ chối.
\end{lemma}
\begin{proof}
$resolveNavigation$ là hàm bộ phận. Do đó hai suy dẫn thành công cho cùng kiểu
receiver, role và kiểu qualifier nhận các bộ bằng nhau. Tra cứu mơ hồ hoặc
không được hỗ trợ là không xác định và không có suy dẫn binding.
\end{proof}

\begin{lemma}[B4 Bảo toàn phạm vi từ vựng và che khuất]
Các biến iterator và let được bind trong cùng phạm vi từ vựng như trong OCL,
kể cả che khuất.
\end{lemma}
\begin{proof}
Các luật binding dùng $\Gamma[x:\tau]$ đúng trong khi suy dẫn thân iterator
hoặc let và khôi phục $\Gamma$ sau đó. Tra cứu gần nhất phân giải $x$ tới
binding mới và bảo toàn mọi $y\neq x$. Quy nạp theo độ sâu lồng nhau chứng minh
mệnh đề; đổi tên alpha bằng tên mới không đổi ngữ nghĩa.
\end{proof}

\begin{lemma}[A1 Bảo toàn kiểu từ Bound sang VA]
Nếu $MM;\Gamma\vdash b:\tau$ và $T_{\mathrm{VA}}(b)=v$ thì
$MM;\Gamma\vdash v:\tau$ trong hệ định kiểu VA.
\end{lemma}
\begin{proof}
Quy nạp cấu trúc trên suy dẫn định kiểu Bound. Các trường hợp self, biến và
literal trực tiếp bảo toàn kiểu đã khai báo. Các trường hợp thuộc tính,
navigation và toán tử kiểu sao chép phần tử siêu mô hình đã phân giải và dùng
luật định kiểu VA tương ứng. Các trường hợp vô hướng và Boolean áp dụng cùng
tiền đề tương thích toán hạng. Các trường hợp iterator, collection, let và if
dùng giả thiết quy nạp dưới cùng mở rộng môi trường; B4 cung cấp kiểu binder và
kỷ luật che khuất. Do đó mỗi constructor được sinh có chính xác kiểu kết quả
nguồn đã dịch, kể cả kiểu phần tử tập hữu hạn.
\end{proof}

\begin{lemma}[A2 Mô phỏng cục bộ từ Bound sang VA]
$T_{\mathrm{VA}}$ là hàm cấu trúc toàn phần trên các term đã binding thành
công, được cho bởi bảng đầy đủ sau, trong đó $V=T_{\mathrm{VA}}$.

\begingroup\small
\begin{longtable}{p{0.48\textwidth}p{0.45\textwidth}}
\toprule Constructor Bound & Kết quả VA\\ \midrule
$BSelf(C)$; $BVar(x,\tau)$; $BLiteral(c,\tau)$ & $Self$; $Var(x)$; $Literal(c)$\\
$BSetLiteral([e_1,\ldots,e_n],Set(\tau))$ & $SetLiteral([V(e_1),\ldots,V(e_n)])$\\
$BAttribute(e,a,\tau)$ & $Attribute(V(e),a)$\\
$BNavigation(e,A,fr,tr,dir,\bar q,Set(E))$ & $Navigation(V(e),A,fr,tr,dir,V(\bar q))$\\
$BAllInstances(C)$; $BNot(e)$ & $AllInstances(C)$; $Not(V(e))$\\
$BBinary(and/or/implies,e_1,e_2)$ & $And/Or/Implies(V(e_1),V(e_2))$\\
$BBinary(op,e_1,e_2)$, so sánh & $Compare(mapCompare(op),V(e_1),V(e_2))$\\
$BBinary(op,e_1,e_2)$, số học & $Arith(mapArith(op),V(e_1),V(e_2))$\\
$BIf(c,t,f,\tau)$ & $If(V(c),V(t),V(f))$\\
$BLet(x,T,e,b,\tau)$ & $Let(x,V(e),V(b))$\\
$BIterator(exists/forAll,S,x,P)$ & $Exists/ForAll(V(S),x,V(P))$\\
$BIterator(select/reject,S,x,P)$ & $Select/Reject(V(S),x,V(P))$\\
$BIterator(collect,S,x,E)$ & $Collect(V(S),x,V(E))$\\
$BIterator(isUnique,S,x,E)$ & $IsUnique(V(S),x,V(E))$\\
$BCollectionOp(includes/excludes,S,[e])$ & $Includes/Excludes(V(S),V(e))$\\
$BCollectionOp(includesAll/excludesAll,S,[T])$ & $IncludesAll/ExcludesAll(V(S),V(T))$\\
$BCollectionOp(union/intersection,S,[T])$ & $Union/Intersection(V(S),V(T))$\\
$BCollectionOp(asSet,S,[])$ & $AsSet(V(S))$\\
$BCollectionOp(size,S,[])$ & $Count(V(S))$\\
$BCollectionOp(isEmpty/notEmpty,S,[])$ & $IsEmpty/NotEmpty(V(S))$\\
$BTypeOp(kindOf,e,C)$ & $TypeKindOf(V(e),C)$\\
$BTypeOp(cast,e,C)$ & $Cast(V(e),C)$\\
\bottomrule
\end{longtable}
\endgroup

$Count$ là phép hạ chuẩn của $size()$ nguồn; nó không phải
$count(element)$ của OCL. $Size$ là từ đồng nghĩa ngữ nghĩa trong VA, được
chuẩn hóa trước các phép viết lại so sánh count. Với mỗi phương trình, đẳng thức
ngữ nghĩa đối tượng của các nút con kéo theo đẳng thức ngữ nghĩa của hai
constructor cục bộ.
\end{lemma}
\begin{proof}
Phân trường hợp trên bảng dịch đầy đủ và khai triển hai ngữ nghĩa cục bộ. Các
trường hợp cơ sở trả về cùng giá trị môi trường hoặc literal. Thuộc tính và
navigation dùng metadata đã phân giải giống nhau, kể cả hướng và thứ tự
qualifier đã đánh giá. Các trường hợp Boolean, so sánh và số học dùng phép toán
vô hướng tương ứng được khai báo độc lập. Iterator và collection áp dụng lập
luận ngoại diên trên tập hữu hạn dưới cùng mở rộng môi trường; $Collect$ là
phép dựng ảnh theo tập hữu hạn. Set literal và hợp/giao dùng cùng
$setJoin$ và phép nhúng chuẩn của M3a ở hai ngữ nghĩa; $asSet$ là đồng nhất.
$IsUnique$ đúng ở cả hai phía khi và chỉ khi ảnh thân không có va chạm theo
đẳng thức có kiểu. $Let$ dùng A3 cho tra cứu mở rộng và $If$ chọn
cùng nhánh theo $\boolval$. Kind-of và cast dùng cùng quan hệ phù hợp. Vì vậy,
kết quả con bằng nhau kéo theo kết quả cục bộ bằng nhau ở mọi dòng. Đây vẫn là
bổ đề mô phỏng cục bộ; không dùng mã hóa đồ thị.
\end{proof}

\begin{lemma}[A3 Bảo toàn tra cứu môi trường]
$T_{\mathrm{VA}}$ bảo toàn tên của $self$, biến tự do, biến iterator và biến
let. Các môi trường tương ứng có cùng miền và giá trị phía đối tượng; mở rộng
stack giống nhau bảo toàn tra cứu gần nhất và che khuất.
\end{lemma}
\begin{proof}
$T_{\mathrm{VA}}$ bảo toàn định danh khai báo gắn với mỗi lần xuất hiện biến.
Định nghĩa tương ứng môi trường theo từng điểm bằng đẳng thức tại mỗi định danh
đó. Ban đầu quan hệ đúng với $self$ và các biến tự do. Nếu nó đúng trước một
binder iterator hoặc let, mở rộng cả hai ánh xạ bằng cùng định danh và giá trị
làm tra cứu mới bằng nhau, đồng thời giữ nguyên mọi tra cứu khác. Khi thoát
phạm vi, cả hai ánh xạ trước đó được khôi phục. Quy nạp theo độ lồng binder do
đó chứng minh đẳng thức tra cứu, kể cả che khuất và đổi tên alpha tránh bắt giữ.
\end{proof}

\begin{definition}[Tập luật chuẩn hóa $R$]
Bộ chuẩn hóa dưới-lên tất định sử dụng:
\begin{align*}
Size(S)&\to Count(S),\\
ForAll(S,x,P)&\to Not(Exists(S,x,Not(P))),\\
NotEmpty(S)&\to Exists(S,z,true),\\
IsEmpty(S)&\to Not(Exists(S,z,true)),\\
Reject(S,x,P)&\to Select(S,x,Not(P)),\\
Implies(A,B)&\to Or(Not(A),B),\\
Xor(A,B)&\to Or(And(A,Not(B)),And(Not(A),B)),\\
Includes(S,y)&\to Exists(S,x,Compare(EQ,Var(x),y)),\\
Excludes(S,y)&\to Not(Exists(S,x,Compare(EQ,Var(x),y))),\\
Compare(GT,Count(N),0)&\to Exists(N,z,true),\\
Compare(EQ,Count(N),0)&\to Not(Exists(N,z,true)),
\end{align*}
trong đó $N$ là một navigation và mọi $x,z$ được đưa vào đều là tên mới. Phép
viết lại tránh bắt giữ và các term được đồng nhất tới tương đương alpha.
\end{definition}

Danh sách này định nghĩa phép biến đổi tham chiếu hình thức $T_{NORM}$; nó không
khẳng định Java phải tạo cùng cây cú pháp. \texttt{OclIrOptimizer} cài một
$T_{OPT}$ hướng graph: \texttt{Implies}/\texttt{Xor} được khai triển trực tiếp,
còn các trường hợp navigation ForAll/rỗng/filter/size có thể trở thành
\texttt{NavigationPredicateCheck} hoặc \texttt{NavigationCountComparison};
collection operation tổng quát có thể được giữ lại cho planner. Vì vậy điều
cần chứng minh là refinement ngữ nghĩa $T_{OPT}(va)\simeq T_{NORM}(va)$, không
phải đẳng thức cú pháp. Test rewrite/differential hiện chỉ là evidence hữu hạn;
proof typed và tránh bắt giữ phổ quát vẫn mở.

\begin{lemma}[N1 Bảo toàn ngữ nghĩa của phép viết lại]
Mỗi phép viết lại chuẩn hóa bảo toàn ngữ nghĩa VA.
\end{lemma}
\begin{proof}
Kiểm tra từng luật đã nêu dưới ngữ nghĩa kiểm tra theo tập hữu hạn. Chân lý phổ
quát là không có phản ví dụ; không rỗng và rỗng lần lượt là tồn tại và không
tồn tại; reject là select theo vị từ kiểm tra bị phủ định; kéo theo là
$\neg A\lor B$; xor đúng chính xác khi một trong hai giá trị
$\boolval(A),\boolval(B)$ là true, đúng bằng
$Or(And(A,Not(B)),And(Not(A),B))$; quan hệ thuộc là đẳng thức tồn tại; loại trừ là phủ định quan
hệ thuộc; $Size$ và $Count$ biểu thị cùng lực lượng hữu hạn; lực lượng
navigation hữu hạn dương hoặc bằng không chính xác khi tồn tại hoặc không tồn
tại witness. Các biến iterator mới tránh bắt giữ. Vì vậy hai vế có cùng ngữ
nghĩa dưới mọi môi trường hợp lệ.
\end{proof}

\begin{lemma}[N2 Bảo toàn kiểu của phép viết lại]
Mỗi phép viết lại chuẩn hóa bảo toàn kiểu tĩnh.
\end{lemma}
\begin{proof}
Xem xét các luật định kiểu cho thấy mọi mẫu nguồn Boolean được thay bằng một
biểu thức Boolean, còn $Reject(S,x,P)$ và $Select(S,x,Not(P))$ có cùng kiểu
phần tử collection. Các biến mới nhận kiểu phần tử nguồn. Do đó mọi luật bảo
toàn kiểu kết quả.
\end{proof}

\begin{longtable}{p{0.13\textwidth}p{0.23\textwidth}p{0.36\textwidth}p{0.20\textwidth}}
\toprule
Luật & Tiền đề & Phương trình ngữ nghĩa & Nghĩa vụ kiểu/độ đo\\
\midrule
N-Size & $S:Set(\tau)$ & $Size(S)=|\finiteSet(S)|=Count(S)$ & $Integer$; $\#Size-1$\\
N-ForAll & $S:Set(\tau)$; $P:Boolean$ & chân lý phổ quát khi và chỉ khi không có witness của $Not(P)$ & $Boolean$; $\#ForAll-1$\\
N-NotEmpty & $S:Set(\tau)$; $x$ mới & không rỗng khi và chỉ khi có witness đúng & $Boolean$; $\#NotEmpty-1$\\
N-IsEmpty & như trên & rỗng khi và chỉ khi không có witness đúng & $Boolean$; $\#IsEmpty-1$\\
N-Reject & Boolean $P$ dưới $x:\tau$ & reject theo $P$ bằng select theo $Not(P)$ & $Set(\tau)$; $\#Reject-1$\\
N-Implies & $A,B:Boolean$ & kéo theo bằng $Or(Not(A),B)$ & $Boolean$; $\#Implies-1$\\
N-Xor & $A,B:Boolean$ & đúng khi chính xác một toán hạng validation-true & $Boolean$; $\#Xor-1$\\
N-Includes & $y$ so sánh được; $x$ mới & quan hệ thuộc khi và chỉ khi tồn tại $\eqval(x,y)$ & $Boolean$; $\#Includes-1$\\
N-Excludes & như trên & loại trừ khi và chỉ khi không có witness đẳng thức & $Boolean$; macro loại cả hai dạng nguồn\\
N-CountPos & $Nav:Set(E)$; $x$ mới & $|Nav|>0$ khi và chỉ khi có witness & $Boolean$; mẫu count $-1$\\
N-CountZero & như trên & $|Nav|=0$ khi và chỉ khi không có witness & $Boolean$; mẫu count $-1$\\
\bottomrule
\end{longtable}

Mọi phương trình đúng trong cả hai diễn giải. Tránh bắt giữ suy ra từ bổ đề thế
chuẩn
$\sem{e[x:=y]}{I}(\rho[y\mapsto v])=\sem{e}{I}(\rho[x\mapsto v])$ với $y$ mới,
được chứng minh theo cấu trúc. Không đích nào nhân đôi một term con
chứa redex tại một bước của chiến lược: phép duyệt dưới-lên chỉ viết lại nút
cha sau khi mọi nút con đã ở dạng chuẩn. Đặc biệt, mặc dù N-Xor sao chép $A$ và
$B$, hai term này không còn chứa constructor được $\mathcal M$ đếm.

\begin{lemma}[N3 Tính dừng của chuẩn hóa]
Chiến lược chuẩn hóa tất định luôn dừng. Cụ thể, định nghĩa:
\[
\mu(e)=
\#ForAll(e)+\#Implies(e)+\#Xor(e)+\#Reject(e)+\#IsEmpty(e)+\#NotEmpty(e)
+\#Includes(e)+\#Excludes(e)+\#CountCompareOnNavigation(e).
\]
Đặt $\mathcal M(e)=(\#Size(e),\mu(e))$ với thứ tự từ điển. Khi đó, với mọi bước
viết lại chiến lược đơn $e\to_{\mathrm{norm}}e'$,
\[
\mathcal M(e')<_{lex}\mathcal M(e).
\]
\end{lemma}
\begin{proof}
Độ đo nhận giá trị trong $\mathbb{N}^2$. $Size\to Count$ làm giảm thành phần
đầu và không bao giờ tái tạo $Size$. Chuẩn hóa vận hành quy định
$Excludes(S,y)$ là macro-rule
$Not(Exists(S,x,Compare(EQ,Var(x),y)))$. Mọi luật chuẩn hóa một bước ngoài phép
chuẩn hóa size đều giữ nguyên thành phần đầu, loại ít nhất một constructor được
$\mu$ đếm và không đưa constructor nào như vậy vào tại gốc. Với N-Xor, các cây
con bị sao chép đã ở dạng chuẩn theo chiến lược dưới-lên, nên việc sao chép
không tăng $\#Size$ hoặc $\mu$. Do đó độ đo từ điển giảm
nghiêm ngặt và không tồn tại chuỗi chuẩn hóa vô hạn.
\end{proof}

\begin{lemma}[N4 Bao phủ dạng chuẩn]
Cho $Redex_R(t)$ nghĩa là $t$ khớp vế trái của một luật đã nêu, và định nghĩa
\[
NF_R(e)\iff \forall t\preceq e:\neg Redex_R(t).
\]
Với mọi biểu thức VA định kiểu đúng $e$, $NF_R(T_{\mathrm{NORM}}(e))$. Phép
viết lại tránh bắt giữ và các term được đồng nhất tới tương đương alpha. Tính
bao phủ là tương đối đối với $R$.
\end{lemma}
\begin{proof}
Dùng quy nạp quy chuẩn trên $\mathcal M(e)$ lồng với quy nạp cấu trúc. Phép
duyệt dưới-lên trước hết trả về các nút con dạng chuẩn. Nếu nút cha được dựng
lại là redex, luật chiến lược duy nhất của nó làm giảm nghiêm ngặt $\mathcal M$
theo N3 và giả thiết quy nạp chuẩn hóa kết quả. Nếu không, mọi term con thỏa
$NF_R$. Vì vậy không redex đã nêu nào còn lại.
\end{proof}

\begin{remark}
Tương ứng thuộc tính và navigation vẫn là R4 và R6. Các bổ đề G nâng những kết
quả biểu diễn nguyên thủy đó qua môi trường, tập hữu hạn, giá trị vô hướng và
chính sách ở mức kiểm tra thay vì lặp lại chúng.
\end{remark}

\begin{lemma}[G1 Bảo toàn mã hóa môi trường]
Nếu môi trường đối tượng $\rho$ và môi trường đồ thị $\eta$ thỏa
$\rho\sim_{\Phi}\eta$, thì mở rộng cả hai môi trường bằng các giá trị tương
ứng bảo toàn quan hệ:
\[
\rho[x\mapsto v]\sim_{\Phi}\eta[x\mapsto \encodeValue(v)].
\]
\end{lemma}
\begin{proof}
Tương ứng môi trường được định nghĩa theo từng điểm. Các binding hiện có giữ
nguyên, còn binding mới tương ứng theo định nghĩa vì môi trường đồ thị lưu
$\encodeValue(v)$. Do đó các môi trường mở rộng thỏa cùng quan hệ theo từng điểm.
\end{proof}

\begin{lemma}[G2 Tính đơn ánh của mã hóa giá trị và đồng cấu tập hữu hạn]
Trên miền giá trị có kiểu được định lý hỗ trợ,
\[
\encodeValue(v_1)=\encodeValue(v_2)\iff v_1=v_2.
\]
Do đó, với các tập hữu hạn $A,B$, $\encodeValue$ bảo toàn và phản ánh quan hệ
thuộc, hợp, tập con, giao, tính rỗng và lực lượng:
\[
\encodeValue(A\cup B)=\encodeValue(A)\cup\encodeValue(B),
\qquad
\encodeValue(A\cap B)=\encodeValue(A)\cap\encodeValue(B),
\]
\[
A\subseteq B \iff \encodeValue(A)\subseteq\encodeValue(B),
\qquad
|A|=|\encodeValue(A)|.
\]
\end{lemma}
\begin{proof}
Tính đơn ánh được chứng minh bằng quy nạp trên cấu trúc giá trị. Thực thể đối
tượng dùng R1/R2, giá trị vô hướng được hỗ trợ dùng mã hóa đồng nhất, còn bottom
và mọi constructor là các loại gắn thẻ rời nhau. Với tập hữu hạn, đẳng thức
ngoại diên và giả thiết quy nạp phản ánh đẳng thức từng phần tử. Ảnh theo từng
phần tử sau đó bảo toàn quan hệ thuộc và hợp; tính đơn ánh phản ánh quan hệ
thuộc, do đó bảo toàn tập con và giao theo cả hai chiều. Nó tạo song ánh giữa
mỗi tập hữu hạn và ảnh của tập, bảo toàn tính rỗng và lực lượng.
\end{proof}

\begin{lemma}[G3 Bảo toàn đẳng thức có kiểu và toán tử vô hướng]
Với mọi giá trị được định lý hỗ trợ,
\[
\eqval(v_1,v_2)\iff
\eqval(\encodeValue(v_1),\encodeValue(v_2)).
\]
Trên các kiểu vô hướng tương thích được hỗ trợ, phép mã hóa còn bảo toàn thứ tự,
số học tại nơi xác định và chân lý kiểm tra:
\[
x=y \iff \encodeValue(x)=\encodeValue(y),
\qquad
x<y \iff \encodeValue(x)<\encodeValue(y),
\]
\[
\encodeValue(x\oplus y)=\encodeValue(x)\oplus\encodeValue(y)
\quad\text{khi }\oplus\text{ được xác định},
\qquad
\boolval(x)=\boolval(\encodeValue(x)).
\]
Nếu $\tau\leq\upsilon$ là phép nâng do $setJoin$ chọn, thì
\[
\encodeValue(\iota^{obj}_{\tau\hookrightarrow\upsilon}(v))
=\iota^{graph}_{\tau\hookrightarrow\upsilon}(\encodeValue(v)),
\]
và phương trình tương tự đúng theo từng phần tử cho $\operatorname{upSet}$.
\end{lemma}
\begin{proof}
Đẳng thức có kiểu suy ra từ G2 và các loại giá trị rời nhau; với thực thể, đó
là định danh đối tượng/nút. Giá trị vô hướng được hỗ trợ dùng cùng giá trị chuẩn
trong miền đối tượng và đồ thị, nên thứ tự tương thích được phản ánh và bảo
toàn. Với phép toán số học xác định ở cả hai phía, áp dụng nó trước hoặc sau mã
hóa đồng nhất cho cùng giá trị vô hướng. Đồng nhất Boolean cũng bảo toàn
$\boolval$. Với phép nâng kiểu, trường hợp đồng nhất và bottom là trực tiếp;
Integer--Real dùng cùng phép đổi chính xác do $ScalarClosed$ bảo đảm; phép nâng
lớp giữ cùng đối tượng/nút nhờ R1--R3. Phương trình cho tập suy ra bằng ngoại
diên từ phương trình phần tử.
\end{proof}

\begin{lemma}[G4 Tương thích với chính sách kiểm tra]
Với các giá trị được định lý hỗ trợ,
\[
\finiteSet(\encodeValue(v))
=\encodeValue(\finiteSet(v)),
\qquad
\asSources(\encodeValue(v))
=\encodeValue(\asSources(v)),
\qquad
\boolval(\encodeValue(v))=\boolval(v).
\]
Đẳng thức đầu chỉ áp dụng tại vị trí collection định kiểu đúng; đẳng thức thứ
hai chỉ áp dụng tại vị trí nguồn navigation định kiểu đúng. Chúng gồm các chính
sách bottom-thành-rỗng đã nêu và không khẳng định lan truyền null/invalid đầy
đủ của OMG OCL.
\end{lemma}
\begin{proof}
Nếu $v=Set(S)$, hai vế của đẳng thức đầu đều là ảnh mã hóa theo từng phần tử của
$S$. Nếu $v=\bot$, hai vế đều rỗng theo chính sách đã nêu. Đây là các trường
hợp vị trí collection duy nhất trong miền định lý. Đẳng thức chân lý đúng vì
true được mã hóa thành true và mọi giá trị được hỗ trợ khác đều không đúng theo
kiểm tra ở cả hai phía.
\end{proof}

\begin{definition}[Giao diện ngữ nghĩa Cypher giới hạn]
$Cypher_{val}$ chỉ chứa alias được sinh, tham số, hằng dành riêng, tra cứu thuộc
tính, biểu thức Boolean/so sánh/số học, \texttt{coalesce}, \texttt{CASE},
\texttt{EXISTS}, \texttt{COUNT}, cùng truy vấn được ghép từ \texttt{MATCH},
\texttt{WHERE}, \texttt{WITH}, \texttt{UNWIND}, \texttt{CALL} tương quan,
\texttt{UNION ALL} và \texttt{RETURN DISTINCT}. Cho $Val_G^C$ là miền giá trị
đồ thị được định lý hỗ trợ và $CyVal_C$ là các giá trị runtime do profile này
sinh ra. Giá trị ngoài $CyVal_C$ không được gán ý nghĩa OCL. Giao diện biểu diễn
có kiểu là
\[
\reprC:Val_G^C\to CyVal_C,\qquad
\absC:CyVal_C\to Val_G^C,
\]
\[
\absC(\reprC(v))=v,\qquad
\absC(CypherNull)=\bot,\qquad
\reprC(\bot)=BOTTOM\_TOKEN.
\]
Một phép trừu tượng hóa bộ phận phụ ánh xạ giá trị Cypher tùy ý thành giá trị đồ
thị được hỗ trợ hoặc $unsupported$; mọi phán đoán hiện thực hóa chỉ lượng hóa
trên các thực thi đóng theo profile đích trong $CyVal_C$. $BOTTOM\_TOKEN$ khác
null được dành riêng nằm ngoài các mã hóa giá trị mô hình và được dùng khi
bottom là phần tử của tập hữu hạn, để phép đếm không loại nó như null. Chính
xác hơn: $\absC\circ\reprC$ là đồng nhất; $\reprC$ đơn ánh;
$BOTTOM\_TOKEN$ không va chạm với giá trị vô hướng, thực thể, khóa lớp,
qualifier payload hay định danh đối tượng; phép tra cứu bảo toàn giá trị đã biểu
diễn; \texttt{DISTINCT} giữ chính xác một lần xuất hiện của token khác null; và
\texttt{COUNT(DISTINCT x)} đếm nó một lần. Trường hợp vô hướng không xác định
có thể dùng $CypherNull$, nhưng bottom ngữ nghĩa được chiếu vào tập phải được
đổi thành token trước. Profile cụ thể cung cấp token dưới dạng tham số dành
riêng \texttt{\$oclBottom} có giá trị map gắn thẻ
\texttt{\{\_\_oclBottom:true\}}. Map không thuộc bất kỳ miền Boolean, số,
String, thực thể, định danh hoặc metadata mô hình được chấp nhận nào. Vì vậy
tính tách biệt theo loại là cấu trúc, không dựa trên việc chọn một chuỗi hiếm.
Về hình thức, $BottomSeparated(MM,M,e)$ yêu cầu adapter không cho phép map gắn
thẻ này xuất hiện như dữ liệu mô hình, khóa hay qualifier payload. Với ánh xạ biến-sang-alias
đơn ánh $\alpha$, định nghĩa
\[
\RowCorr(\alpha,\eta,r)\iff
\forall x\in\operatorname{dom}(\eta):\absC(r(\alpha(x)))=\eta(x).
\]
Với bag dòng $B$,
$\projectSet(a,B)=\{\absC(r(a))\mid r\text{ xuất hiện trong }B\}$.
Ngữ nghĩa biểu thức và dòng được ký hiệu bởi
$EvalExpr_{G,\sigma}(c,r)$ và $EvalRows_{G,\sigma}(Q,B)$.
\end{definition}

\begin{assumption}[Hợp đồng biểu diễn đích BR1--BR7]
Mọi thực thi thuộc profile đích thỏa các luật sau.
\begin{enumerate}[label=(BR\arabic*)]
  \item $\absC(\reprC(v))=v$ với mọi $v\in Val_G^C$.
  \item $\reprC$ đơn ánh trên từng miền giá trị có kiểu.
  \item $\absC(CypherNull)=\bot$ ở vị trí vô hướng và
  $\reprC(\bot)=BOTTOM\_TOKEN$ khi bottom được vật chất hóa như một phần tử tập.
  \item $BOTTOM\_TOKEN$ khác null và không va chạm với mọi giá trị mô hình,
  định danh, khóa hoặc qualifier payload có thể đạt tới.
  \item Tra cứu tham số, thuộc tính và alias của một giá trị đã biểu diễn bảo
  toàn giá trị đó dưới $\absC$.
  \item \texttt{DISTINCT} trên một cột đã qua $setOut$ giữ đúng một đại diện
  cho mỗi giá trị ngữ nghĩa, kể cả đúng một đại diện cho bottom.
  \item \texttt{COUNT(DISTINCT x)} trên cột đã qua $setOut$ đếm lực lượng tập
  ngữ nghĩa; đặc biệt token bottom khác null được đếm đúng một lần.
\end{enumerate}
\end{assumption}

\begin{definition}[Kế hoạch vô hướng và bảo toàn alias đang sống]
Với tập hữu hạn alias đang sống $K$, viết
\[
Preserves(K,r,r')\iff
\forall a\in K:\ a\in dom(r)\cap dom(r')\land
\absC(r'(a))=\absC(r(a)).
\]
Đặt $ExecPrelude(L,r)=EvalRows_{G,\sigma}(Seq(L),\{r\})$. Định nghĩa
\[
WFScalar(K,ScalarPlan(L,c))
\]
khi và chỉ khi mọi dòng vào được chấp nhận sinh chính xác một dòng $r'$, bảo
toàn $K$ và $FV(c)\subseteq dom(r')$. Tập đang sống chứa $range(\alpha)$ và mọi
alias mới cần cho phần tiếp diễn. Vì vậy mỗi \texttt{WITH} được sinh chiếu đồng
nhất từng alias đang sống bên cạnh binding mới của nó.
\end{definition}

\begin{definition}[Các phán đoán hiện thực hóa]
Với mọi môi trường và dòng tương ứng, hiện thực hóa vô hướng, vị từ và tập hữu
hạn lần lượt có nghĩa là:
\[
\alpha,\sigma\vdash c\ \mathsf{realizes}_{scalar}\ v:\tau
\iff
\absC(EvalExpr_{G,\sigma}(c,r))=\sem{v}{graph}^{\tau}(G,\eta),
\]
và hiện thực hóa kế hoạch vô hướng là
\[
\alpha,\sigma;K\vdash ScalarPlan(L,c)\ \mathsf{realizes}_{scalar}\ v:\tau
\]
khi và chỉ khi $WFScalar(K,ScalarPlan(L,c))$ và, với dòng duy nhất
$r'\in ExecPrelude(L,r)$,
\[
\RowCorr(\alpha,\eta,r')\land
\absC(EvalExpr_{G,\sigma}(c,r'))=\sem{v}{graph}^{\tau}(G,\eta).
\]
\[
\alpha,\sigma\vdash c\ \mathsf{realizes}_{pred}\ v
\iff
\boolval(\absC(EvalExpr_{G,\sigma}(c,r)))
=\boolval(\sem{v}{graph}^{Boolean}(G,\eta)),
\]
\[
\alpha,\sigma\vdash Q\Rightarrow a\ \mathsf{realizes}_{set}\ v:Set(\tau)
\iff
\projectSet(a,EvalRows_{G,\sigma}(Q,\{r\}))
=\finiteSet(\sem{v}{graph}^{Set(\tau)}(G,\eta)),
\]
trong đó lượng hóa phổ quát trên $\eta,r$ thỏa $\RowCorr(\alpha,\eta,r)$.
Phán đoán vị từ của kế hoạch vô hướng thay đẳng thức cuối bằng đẳng thức dưới
$\boolval$. Hiện thực hóa chỉ bằng biểu thức là trường hợp riêng $L=[]$.
Với context $C$,
\[
\sigma\vdash q\ \mathsf{realizes}_{inv}(C,nva)
\iff
\returnedIds(q,G,\sigma)=
\{\id(o)\mid o\in\Obj(M),\ classOf(o)\mathrel{conformsTo}C,
\ \boolval(\sem{nva}{graph}^{Boolean}
   (G,self\mapsto\node(o)))=false\}.
\]
Phán đoán tập loại bỏ bội số dòng; phán đoán invariant so sánh các định danh,
không so sánh bảng kết quả hoặc nút được trả về.
\end{definition}

\begin{assumption}[Các tiên đề hành vi của $Cypher_{val}$]
Ngữ nghĩa ngoại diên được chọn cho AST thô thỏa chín luật sau.
\begin{enumerate}[label=(CY\arabic*)]
  \item Alias, tham số, thuộc tính, Boolean, so sánh, số học,
  \texttt{CASE} và \texttt{coalesce} được đánh giá theo profile trên các giá
  trị được hỗ trợ; thuộc tính vắng mặt cho $CypherNull$ và $truth$ chỉ cho true
  đối với Boolean true.
  \item \texttt{MATCH} liệt kê đúng các phép gán pattern thỏa đồ thị, với bội
  số bag nhưng không tạo phép gán không khớp.
  \item \texttt{WHERE} giữ đúng các dòng có guard Cypher true; guard do
  renderer sinh là toàn phần nhờ $truth$ hoặc $isBottom$.
  \item \texttt{EXISTS\{Q\}} đúng khi và chỉ khi đánh giá tương quan của $Q$
  sinh ít nhất một dòng.
  \item \texttt{COUNT(DISTINCT x)} tuân theo BR6--BR7 và vì thế đếm lực lượng
  tập sau $setOut$.
  \item \texttt{WITH}/\texttt{RETURN} chiếu các biểu thức đã nêu;
  \texttt{DISTINCT} đồng nhất chính xác các dòng có cùng tuple chiếu theo BR2
  và BR6.
  \item \texttt{UNWIND} sinh một dòng cho mỗi vị trí của danh sách hữu hạn;
  phép chiếu phân biệt phía sau khôi phục ngữ nghĩa tập.
  \item \texttt{CALL} tương quan chỉ nhìn thấy các alias được nhập tường minh,
  đánh giá subquery cho từng dòng ngoài và nối kết quả với các binding nhập
  không đổi.
  \item \texttt{UNION ALL} nối hai bag dòng. Khi hai nhánh có guard toàn phần
  bù nhau và cùng schema, chính xác một nhánh đóng góp; phép chiếu phân biệt
  ngoài cùng khôi phục tập ngữ nghĩa.
\end{enumerate}
\end{assumption}

\begin{lemma}[C1 Bảo toàn alias và phép chiếu đang sống]
Nếu $\RowCorr(\alpha,\eta,r)$, alias $a$ là mới và
$w=\encodeValue(v)$ thì
\[
\RowCorr(\alpha[x\mapsto a],\eta[x\mapsto w],r[a\mapsto\reprC(w)]).
\]
Khôi phục các ánh xạ ngoài sau một subquery bảo toàn các binding bị che khuất.
Ngoài ra, mọi \texttt{WITH} được sinh thỏa $Preserves(K,r,r')$ với tập đầu ra
đang sống $K$ đã tính; vì vậy nó không thể loại \texttt{self}, alias iterator
hay alias ngoài tương quan mà phần tiếp diễn cần.
\end{lemma}
\begin{proof}
Đẳng thức mới đúng tại $x$ theo phép dựng và mọi đẳng thức cũ không đổi nhờ
tính mới. Constructor \texttt{withBind} của phép khai triển toàn phần chiếu
đồng nhất mọi alias trong $K$ và chỉ thêm binding mới, qua đó chứng minh bảo
toàn. Khôi phục stack cho phép che khuất từ vựng.
\end{proof}

\begin{lemma}[C2 Tính đúng đắn của tham số]
Nếu tham số $p$ biểu diễn giá trị được hỗ trợ $s$ thì
$\sigma(p)=\reprC(\encodeValue(s))$ và
$\absC(EvalExpr_{G,\sigma}(\texttt{\$p},r))=\encodeValue(s)$ với mọi dòng $r$.
\end{lemma}
\begin{proof}
Cấp phát tham số và chèn vào $\sigma$ xảy ra trong cùng bước renderer; tên mới
ngăn ghi đè. Tra cứu tham số sau đó dùng tiên đề $Cypher_{val}$ đã nêu.
\end{proof}

\begin{lemma}[C3 Hiện thực hóa truy cập đồ thị nguyên thủy]
Dưới $\RowCorr$, các alias được sinh, tra cứu thuộc tính, pattern quan hệ thuộc
kiểu, accessor thuộc tính chuẩn, pattern all-instances và pattern navigation
nhị phân có hướng thỏa phán đoán hiện thực hóa áp dụng cho $Self$, $Var$,
$Literal$, $Attribute$, $AllInstances$ và $Navigation$. Navigation gồm đánh
giá và khớp danh sách qualifier có thứ tự. Các trường hợp nguyên thủy thực thể
ở đây giả thiết receiver khác bottom đã được biểu diễn; C3a cung cấp phép nâng
an toàn với bottom.
\end{lemma}
\begin{proof}
Alias dùng tương ứng dòng và literal dùng C2. Thuộc tính dùng R4, pattern kiểu
dùng R3/R7, còn pattern navigation ràng buộc kiểu quan hệ,
exact \texttt{associationKey}, role, hướng và danh sách qualifier từ
\texttt{sourceQualifiers}/\texttt{targetQualifiers}; R5/R6
cho đẳng thức của tập ngữ nghĩa được chiếu.
\end{proof}

\begin{lemma}[C3a Receiver thực thể an toàn với bottom]
Cho kế hoạch vô hướng $P$ hiện thực hóa receiver thực thể tĩnh có thể là bottom,
và cho $F$ là bộ dựng thân thực thể đơn hiện thực hóa phép nguyên thủy $f$ cho
mọi thực thể đã biểu diễn trong khi bảo toàn tập đang sống $K$. Nếu $b$ biểu
diễn $f(\bot)$ thì $withEntityReceiver(P,K,b,F)$ là kế hoạch vô hướng hợp
dạng, hiện thực hóa $f$ trên toàn bộ miền receiver được chấp nhận.
\end{lemma}
\begin{proof}
Với bottom, $isBottom$ toàn phần chỉ bật nhánh không có pattern. Nếu không,
định kiểu và tính đóng profile cho một nút và chỉ bật $F$. Tiên đề lọc ngăn các
clause sau guard false nhận dòng; ngữ nghĩa call tương quan bảo toàn các binding
đã nhập và trả về kết quả đơn được bật. Do đó chính xác một dòng được sinh và
C1 cho $WFScalar$. Thuộc tính, kind-of và cast lần lượt cụ thể hóa kết quả bằng
R4, R3 và các kết quả bottom là null, false, null.
\end{proof}

\begin{lemma}[C4 Hợp thành vị từ và lượng từ tồn tại]
Nếu các fragment con thỏa phán đoán hiện thực hóa, biểu thức Boolean và so sánh
được sinh bảo toàn $\boolval$; \texttt{WHERE} giữ chính xác các dòng đúng theo
kiểm tra; và \texttt{EXISTS} đúng theo kiểm tra chính xác khi truy vấn witness
đã hiện thực hóa không rỗng. Do đó các phán đoán hợp thành cho $Not$, $And$,
$Or$, so sánh, $If$, $Let$, $Exists$ và các dạng không-có-phản-ví-dụ đã chuẩn hóa.
\end{lemma}
\begin{proof}
Áp dụng các giả thiết quy nạp cho biểu thức con và các tiên đề tường minh về
chân lý, lọc, tồn tại, \texttt{CASE}, \texttt{coalesce} của $Cypher_{val}$.
C1 cung cấp các mở rộng iterator và let.
\end{proof}

\begin{lemma}[C5 Hợp thành tập hữu hạn và lực lượng]
Khi nguồn và thân đã được hiện thực hóa, các fragment lọc, ảnh, quan hệ thuộc,
tập con, rời nhau, hợp, giao, chuyển thành tập, tính duy nhất, tính rỗng và lực lượng được sinh thỏa các phán đoán hiện thực
hóa tương ứng. Phép chiếu ngữ nghĩa là tập, và lực lượng dùng \texttt{DISTINCT}
khi nhiều dòng có thể biểu thị một phần tử VA. Bottom ngữ nghĩa trong tập được
chiếu thành $BOTTOM\_TOKEN$ khác null.
\end{lemma}
\begin{proof}
Phép chiếu tập ngoại diên loại bội số dòng. Phép lọc dùng C4; phép dựng ảnh dùng
C1 trên mỗi phần tử nguồn. M3a và G3 cho phép coercion mọi phần tử tới kiểu
$setJoin$ chung trước $setOut$; BR1--BR6 làm $setEq$ phản ánh đúng
$\eqval$ và làm phép chiếu phân biệt đúng với tập ngữ nghĩa. Vì thế Set literal
là ảnh hữu hạn của danh sách phần tử, union là hợp ngoại diên, intersection là
tập các phần tử phía trái có witness bằng nó ở phía phải, còn $asSet$ là đồng
nhất. Với $isUnique$, nguồn đã phân biệt cung cấp đúng một dòng cho mỗi phần tử;
BR7 và CY5 làm đẳng thức giữa số dòng nguồn và số giá trị thân phân biệt tương
đương với tính đơn ánh của ảnh thân. Cùng luật aggregation cho lực lượng, còn
membership, tập con và rời nhau theo các định nghĩa witness/counterexample.
\end{proof}

\begin{lemma}[C6 Hiện thực hóa lớp bọc invariant]
Nếu $\alpha(self)=a_{self}$ và
$\alpha,\sigma\vdash c\ \mathsf{realizes}_{pred}\ nva$ thì lớp bọc gồm
\texttt{MATCH} context, \texttt{WHERE} vi phạm và
\texttt{RETURN DISTINCT} $a_{self}$\texttt{.use\_id} được sinh thỏa
$\sigma\vdash q\ \mathsf{realizes}_{inv}(C,nva)$.
\end{lemma}
\begin{proof}
R3 cho phép liệt kê context chính xác; C4 cho bộ lọc không-đúng-theo-kiểm-tra
chính xác; R2 cho định danh ổn định; và phép chiếu phân biệt sinh tập trong
phán đoán hiện thực hóa invariant.
\end{proof}

\subsection{Các luật $T_{\mathrm{TEXT}}$ định hướng cú pháp}

Renderer được định nghĩa bởi các phán đoán đệ quy tương hỗ
\[
\alpha;\sigma\vdash v\Rightarrow_E c:\tau;\sigma',
\qquad
\alpha;\sigma\vdash v\Rightarrow_S(Q,a):Set(\tau);\sigma'.
\]
Trạng thái tham số tăng đơn điệu và mọi tên được đưa vào đều là tên mới. Đích
được chọn là profile Neo4j Cypher 5 $CYPHER5_{val}$. Driver 5.21.0 không quyết
định ngữ nghĩa server: thực nghiệm phải ghi lại \texttt{dbms.components()} và
chạy ma trận riêng cho CY1--CY9: representation/lookup, pattern và duplicate
paths, filter, existential, count/distinct, projection, scalar/control truth,
unwind và correlated call/union. Bằng chứng không chuẩn tắc ngày 2026-08-01 ghi
nhận 9/9 pass trên Neo4j Kernel 2026.06.0 Enterprise với thành phần Cypher 5 và
cơ sở dữ liệu \texttt{demo}, dùng Java driver dependency 5.21.0. Manifest kiểm
tra bằng máy ghim profile và hash renderer/encoding/probe/harness. Điều này chỉ
xác nhận các đại diện đã chọn trên runtime đó; nó không thay chứng minh hiện
thực hóa và không bao phủ phiên bản server khác. Các phép khai triển nguyên
thủy là:
\[
isBottom(c)=(c\ IS\ NULL\lor c=BOTTOM\_TOKEN),\quad
truth(c)=coalesce(c=true,false),
\]
\[
setOut(c)=CASE\ WHEN\ c\ IS\ NULL\ THEN\ BOTTOM\_TOKEN\ ELSE\ c\ END.
\]
Mọi phép toán tập trước hết dùng
\[
setEq^{\tau}(c,d)=semEq^{\tau}(setOut(c),setOut(d)).
\]
Với phép nâng $\tau\leq\upsilon$ do $setJoin$ chọn, combinator
$coerce_C^{\tau\hookrightarrow\upsilon}(c)$ hiện thực hóa
$\iota^{graph}_{\tau\hookrightarrow\upsilon}$: đồng nhất và nâng lớp giữ nguyên
giá trị biểu diễn, còn Integer--Real dùng phép đổi chính xác của profile. Mọi
phần tử Set literal và mọi toán hạng của hợp/giao được coercion tới $\upsilon$
trước khi áp dụng $setOut$, $setEq^{\upsilon}$ hoặc \texttt{DISTINCT}.
Do binder đã cố định $\tau$ bằng các luật so sánh hoặc $setJoin$, $setEq$ là
đẳng thức có kiểu và toàn phần. Token hóa xảy ra trước membership,
\texttt{DISTINCT}, hợp, giao, tính duy nhất và đếm; vì vậy hai bottom được đồng
nhất, còn bottom không thể bị Cypher loại như null hoặc bằng một giá trị mô hình.

Với mọi phép toán vô hướng tiêu thụ thực thể, phép khai triển được chứng nhận
dùng $withEntityReceiver(P,K,b,F)$. Nó vật chất hóa kế hoạch receiver $P$ đúng
một lần thành alias mới $recv$, rồi sinh call tương quan chứa hai nhánh
\texttt{UNION ALL}. Nhánh đầu áp dụng $WHERE\ isBottom(recv)$ và trả về giá trị
trường hợp bottom $b$ mà không có pattern đồ thị. Nhánh thứ hai áp dụng
$WHERE\ NOT\ isBottom(recv)$ trước khi thực thi bộ dựng clause chỉ dành cho
thực thể $F(recv)$. Các guard là Boolean toàn phần bù nhau. Định kiểu tĩnh và
tính đóng profile đích kéo theo receiver khác bottom đi vào nhánh hai là một
nút đồ thị đã biểu diễn. Mỗi nhánh trả về cùng alias đầu ra mới và chính xác
một nhánh được bật đóng góp một dòng; mọi alias trong $K$ được nhập và bảo
toàn. Vì vậy phép dựng an toàn với bottom mà không giả thiết \texttt{CASE} dừng
ngắn mạch một pattern đồ thị.

\begingroup\scriptsize
\begin{longtable}{p{0.15\textwidth}p{0.45\textwidth}p{0.33\textwidth}}
\toprule Combinator & Khai triển cụ thể & Tiền đề và chính sách bottom\\ \midrule
$semEq^\tau(c,d)$ & \texttt{CASE}: cả hai bottom $\mapsto true$; đúng một bottom $\mapsto false$; nếu không $c=d$ & đẳng thức có kiểu/kiểu số chung\\
$setEq^\tau(c,d)$ & $semEq^\tau(setOut(c),setOut(d))$ & token hóa trước so sánh phần tử tập\\
$coerce_C^{\tau\hookrightarrow\upsilon}(c)$ & đồng nhất/nâng lớp giữ biểu diễn; Integer--Real đổi chính xác & M3a, G3 và $ScalarClosed$\\
$semOrd(op,c,d)$ & toán hạng bottom $\mapsto null$; nếu không dùng $c\ op\ d$ native & các kiểu vô hướng có thứ tự được chấp nhận\\
$semArith(op,c,d)$ & toán hạng bottom $\mapsto null$; nếu không dùng $c\ op\ d$ native & miền chung biểu diễn được; mẫu khác không\\
$attr_C(P,a)$ & $withEntityReceiver(P,K,null,F_{attr})$; chỉ nhánh thực thể chứa pattern thuộc tính optional & bottom không bao giờ ở vị trí nút; giá trị thiếu ánh xạ thành bottom\\
$qualList_C(\bar c)$ & null nếu có thành phần bottom; nếu không là danh sách có thứ tự $[c_1,\ldots,c_k]$ & qualifier nguyên thủy, số ngôi cố định, tuần tự hóa đơn ánh\\
$kind_C(P,C)$ & $withEntityReceiver(P,K,false,F_{kind})$ & nhánh bottom không có pattern; nhánh thực thể dùng quan hệ phù hợp đã vật chất hóa\\
$cast_C(P,C)$ & $withEntityReceiver(P,K,null,F_{cast})$ & receiver được đánh giá một lần; bottom/cast thất bại ánh xạ thành bottom\\
navigation thuận & \texttt{MATCH (s)-[r]->(t) WHERE type(r) STARTS WITH 'Link'} cộng các vị từ metadata chuẩn & R5/R6; chiếu đích phân biệt\\
navigation ngược & \texttt{MATCH (t)-[r]->(s)} với vị từ role đảo và cùng kiểm tra kiểu & R5/R6; chiếu đích phân biệt\\
$Count(S)$ & \texttt{COUNT \{Q RETURN DISTINCT setOut(s)\}} & probe count-subquery; token được đếm một lần\\
If có giá trị tập & \texttt{CALL} tương quan có guard, \texttt{UNION ALL}, phép chiếu ngoài phân biệt & schema nhánh một cột bằng nhau; guard toàn phần bù nhau\\
\bottomrule
\end{longtable}
\endgroup

Ở đây \texttt{Link*} chỉ là ký hiệu meta; Cypher thô dùng biến quan hệ không
giới hạn và vị từ \texttt{type(r) STARTS WITH 'Link'} đã hiển thị. Mỗi dòng biểu
thức chứng minh $\absC(EvalExpr(expansion,row))$ bằng ngữ nghĩa đồ thị cần có;
mỗi dòng tập chứng minh đẳng thức qua $\projectSet$. Các chứng minh là phân
tích trường hợp cục bộ dùng luật biểu diễn, hợp đồng vô hướng, R3--R6 và C5.

\begingroup\scriptsize
\begin{longtable}{p{0.13\textwidth}p{0.34\textwidth}p{0.46\textwidth}}
\toprule Luật & Tiền đề & Biểu thức được sinh\\ \midrule
TXT-Self/Var & $\alpha(x)=a$ & $a$\\
TXT-Lit & $p$ mới, $\sigma'(p)=\reprC(c)$ & \texttt{\$p}\\
TXT-Attr & $e\Rightarrow_E P$ & $attr_C(P,a)$ qua $withEntityReceiver$\\
TXT-Not & $e\Rightarrow_E c$ & $NOT\ truth(c)$\\
TXT-And/Or & $e_i\Rightarrow_E c_i$ & $truth(c_1)\ AND/OR\ truth(c_2)$\\
TXT-Implies & như trên & $NOT\ truth(c_1)\ OR\ truth(c_2)$\\
TXT-Compare & $e_i\Rightarrow_E c_i$ & $semEq/semOrd(op,c_1,c_2)$; phủ định đẳng thức cho NEQ\\
TXT-Arith & như trên & $semArith(op,c_1,c_2)$\\
TXT-If-E & điều kiện và các nhánh dịch thành biểu thức & $CASE\ WHEN\ truth(c)\ THEN\ c_t\ ELSE\ c_f\ END$\\
TXT-Let-E & $body[x:=init]\Rightarrow_E c$ & $c$\\
TXT-Exists & $S\Rightarrow_S(Q,s)$; $P\Rightarrow_E c$ dưới $\alpha[x\mapsto s]$ & $EXISTS\{Q\ WITH\ s\ WHERE\ truth(c)\ RETURN\ 1\}$\\
TXT-ForAll & như trên & $NOT\ EXISTS\{Q\ WITH\ s\ WHERE\ NOT\ truth(c)\ RETURN\ 1\}$\\
TXT-Incl/Excl & $S\Rightarrow_S(Q,s)$; $y\Rightarrow_E c_y$ & tồn tại/không tồn tại dòng có $setEq(s,c_y)$\\
TXT-InclAll & $S\Rightarrow_S(Q_s,s)$; $T\Rightarrow_S(Q_t,t)$ & không có $t$ mà không tồn tại $s$ với $setEq(s,t)$\\
TXT-ExclAll & như trên & không có $t$ mà tồn tại $s$ với $setEq(s,t)$\\
TXT-Count & $S\Rightarrow_S(Q,s)$ & $COUNT\{Q\ RETURN\ DISTINCT\ setOut(s)\}$\\
TXT-Empty & $S\Rightarrow_S(Q,s)$ & $NOT\ EXISTS\{Q\ RETURN\ 1\}$\\
TXT-NotEmpty & như trên & $EXISTS\{Q\ RETURN\ 1\}$\\
TXT-IsUnique & $S\Rightarrow_S(Q,s)$; $E:\sigma\Rightarrow_E c$ dưới $x\mapsto s$ & so sánh số phần tử nguồn với $COUNT(DISTINCT\ setOut(c))$ sau coercion có kiểu\\
TXT-Kind/Cast & $e\Rightarrow_E P$ & $kind_C/cast_C(P,C)$ qua $withEntityReceiver$\\
\bottomrule
\end{longtable}

\begin{longtable}{p{0.13\textwidth}p{0.34\textwidth}p{0.46\textwidth}}
\toprule Luật & Tiền đề & Truy vấn tập được sinh\\ \midrule
TXT-AllInst & tham số khóa lớp; $n$ mới & khớp $n$ bằng $ObjectInstanceOf$ đã vật chất hóa; trả về phân biệt $setOut(n)$\\
TXT-SetLit & $e_i:\tau_i\Rightarrow_E c_i$; $setJoin(\bar\tau)=\upsilon$ & coercion từng $c_i$ tới $\upsilon$, $setOut$, rồi loại trùng bằng $setEq^\upsilon$\\
TXT-Nav & các dòng nguồn; danh sách qualifier đã dịch; đích $t$ mới & pattern Link có hướng với liên kết, role, $qualList_C$; trả về phân biệt $setOut(t)$\\
TXT-Select & $S\Rightarrow_S(Q,s)$; vị từ dưới $x\mapsto s$ & $Q\ WHERE\ truth(c)\ RETURN\ DISTINCT\ setOut(s)$\\
TXT-Reject & như trên & $Q\ WHERE\ NOT\ truth(c)\ RETURN\ DISTINCT\ setOut(s)$\\
TXT-Collect & các dòng nguồn; biểu thức thân không phải tập $c$ & $Q\ WITH\ setOut(c)\ AS\ z\ RETURN\ DISTINCT\ z$\\
TXT-Union/Intersection & hai kế hoạch tập; $setJoin(\tau,\sigma)=\upsilon$ & coercion hai phía tới $\upsilon$, hợp hoặc giao bằng $setEq^\upsilon$, rồi chiếu phân biệt\\
TXT-AsSet & một kế hoạch tập & chiếu lại cùng phần tử với $DISTINCT$\\
TXT-If-S & điều kiện; hai nhánh tập & call tương quan có guard bù nhau, $UNION\ ALL$, phép chiếu ngoài phân biệt\\
TXT-Let-S & $body[x:=init]\Rightarrow_S(Q,a)$ & $(Q,a)$\\
\bottomrule
\end{longtable}
\endgroup

Với biểu thức thực thể, phép nâng dòng nguồn sinh một dòng có guard; với tập, nó
dùng suy dẫn $\Rightarrow_S$. Vì vậy nó hiện thực hóa $\asSources$. Hai luật
let đổi tên alpha mọi binder trong thân xuất hiện tự do trong initializer trước
khi thế tránh bắt giữ. Định kiểu, biến tự do và che khuất từ vựng được bảo toàn,
và
\[
|body[x:=init]|\leq |body|+occ_{free}(x,body)\,|init|.
\]
Đệ quy dừng theo độ sâu lồng let cực đại; định lý không tuyên bố cận đa thức cho
kích thước truy vấn. Với If có giá trị tập, cả hai nhánh nhập tường minh cùng
các alias ngoài tự do, chiếu một alias có kiểu chung và dùng các guard toàn
phần bù nhau. Do đó chính xác một nhánh đóng góp cho mỗi dòng vào trước phép
chiếu ngoài phân biệt.

\begin{lemma}[TXT1--TXT4 Các siêu thuộc tính của phép dịch]
Với phép sinh tên mới tất định, mọi suy dẫn dịch thỏa bốn kết luận sau.
\begin{enumerate}[label=(TXT\arabic*)]
  \item Kết quả được xác định duy nhất tới đổi tên alpha của alias mới.
  \item $\Rightarrow_E$ và $\Rightarrow_S$ bảo toàn loại cùng kiểu kết quả.
  \item Mọi constructor đầu ra thuộc đúng mảnh $Cypher_{val}$ đã chọn.
  \item Mọi kế hoạch đầu ra thỏa phán đoán hiện thực hóa vô hướng, vị từ hoặc
  tập tương ứng với constructor đầu vào.
\end{enumerate}
\end{lemma}
\begin{proof}
Quy nạp đồng thời trên các suy dẫn $\Rightarrow_E$, $\Rightarrow_S$ và suy dẫn
nguồn navigation. Với TXT1, constructor và loại tĩnh của đầu vào chọn đúng một
luật; mọi lời gọi đệ quy duy nhất theo giả thiết quy nạp, còn bộ cấp phát tên và
tham số là tất định. Hai lần chạy chỉ có thể chọn các atom mới khác tên, nên có
một song ánh alpha giữ nguyên toàn bộ constructor, import và tham chiếu.

Với TXT2, kiểm tra từng dòng bảng dịch. Literal, alias và toán tử vô hướng giữ
kiểu đã suy dẫn; accessor dùng kiểu kết quả đã phân giải; kế hoạch tập xuất đúng
một cột có kiểu phần tử nguồn hoặc kiểu $setJoin$; iterator mở rộng môi trường
bằng đúng kiểu phần tử. Vì không có luật fallback, không thể sinh kế hoạch sai
loại từ một constructor được chấp nhận.

Với TXT3, mỗi vế phải chỉ dùng constructor raw-AST đã liệt kê. Phép hợp thành
chỉ lồng các constructor đó; helper cấp phát alias/tham số chứ không chèn text
đích. Do đó giả thiết quy nạp kéo theo toàn bộ đầu ra vẫn thuộc
$Cypher_{val}$.

Với TXT4, các luật cơ sở dùng C2/C3. Hợp thành biểu thức dùng C4, mở rộng môi
trường iterator dùng C1, còn phép chiếu tập và lực lượng dùng C5. $setOut$ bảo
toàn bottom ngữ nghĩa dưới dạng token khác null. Các nhánh điều kiện có guard
bù nhau dưới $truth$; CY8--CY9 kết hợp chính xác nhánh được chọn. Các trường hợp
D1--D11 bên dưới cung cấp các suy dẫn hiện thực hóa không trực tiếp, nên quy nạp
bao phủ mọi constructor.
\end{proof}

\subsection{Cypher AST thô và phép khai triển toàn phần}

Các ký hiệu $Q_s,Q_t,Q_f$ trong các luật trước đặt tên cho kết quả suy dẫn đệ
quy đã hoàn tất. Chúng không phải lỗ trống văn bản và không phải constructor
của ngôn ngữ đích. Đích được chứng nhận là AST thô có tham số:
\begin{verbatim}
CExpr  ::= Alias | Param | Null | Bool | Int | Property | List
         | Unary | Binary | Case | Function | ListComp
         | ExistsExpr(CQuery) | CountExpr(CQuery)
Pattern ::= Node(alias,label,properties)
          | Rel(node,alias,direction,optionalType,node)
Clause ::= Match(optional,patterns) | Where(expr) | Unwind(expr,alias)
         | With(distinct,projections) | Return(distinct,projections)
         | Call(importedAliases,query)
CQuery ::= Seq(clauses) | UnionAll(CQuery,CQuery)
\end{verbatim}
Định danh, khóa thuộc tính, nhãn, kiểu quan hệ, toán tử và hàm là các atom từ
những enumeration hữu hạn được chứng nhận. Giá trị mô hình chỉ xuất hiện trong
nút tham số. Vì hiện thực hóa vô hướng có thể cần các clause đứng trước, bộ
khai triển toàn phần dùng các kiểu dữ liệu kiến tạo đóng
\[
 ScalarPlan=([Clause],CExpr),\qquad SetPlan=(CQuery,exportedAlias).
\]
Prelude của kế hoạch vô hướng giữ mọi alias tự do của kết quả trong phạm vi;
kế hoạch tập trả về chính xác schema một cột được xuất. Đây là các record kiến
tạo của ngôn ngữ host, không phải cú pháp đích bổ sung.

Cho $CQExpr$, $CQSet$, $CQSource=CQSourceExpr(CQExpr)+CQSourceSet(CQSet)$ và
$CQInv$ là các constructor query-model có kiểu tương ứng một-một với các dòng
TXT. Định nghĩa bằng khớp constructor:
\[
\begin{array}{rcl}
Expand_E&:&CQExpr\times AliasEnv\times ParamState\times LiveSet
          \longrightarrow ScalarPlan\times ParamState,\\
Expand_S&:&CQSet\times AliasEnv\times ParamState
          \longrightarrow SetPlan\times ParamState,\\
Expand_{Src}&:&CQSource\times AliasEnv\times ParamState
          \longrightarrow SetPlan\times ParamState,\\
Expand_I&:&CQInv\times ParamState
          \longrightarrow CQuery\times ParamState.
\end{array}
\]
Mọi lời gọi dùng tập đang sống chứa $range(\alpha)$. Các helper toàn phần sau
loại bỏ các phép prelude/import từng chỉ mang tính sơ đồ:
\[
withBind(K,e,a)=With(false,[x\mapsto x]_{x\in K}\cup[e\mapsto a]),
\]
\[
bindScalar(ScalarPlan(L,e),a,K)
=ScalarPlan(L\mathbin{+\!+}[withBind(K,e,a)],Alias(a)).
\]
Helper toàn phần $prefixSet$ bọc kết quả $Seq$ hoặc $UnionAll$ tùy ý trong một
$Call$ tương quan, thực thi các clause tiền tố và trả về phân biệt cùng alias đã
xuất. Helper $materialize$ đánh giá kế hoạch vô hướng một lần, bind kết quả vào
alias mới bằng $withBind$, và các clause sau chỉ tham chiếu alias đó. Do đó các
prelude qualifier chạy theo thứ tự khai báo, prelude điều kiện đứng trước guard
nhánh và $Cast(P,C)$ đánh giá $P$ chính xác một lần. Call tương quan nhập tường
minh các alias tự do đã tính; nhánh điều kiện vô hướng và tập trả về cùng schema
một cột.

Đối với các constructor tập, $materializeList([P_1,\ldots,P_n],K)$ là phép gấp
trái của $materialize$, sinh các alias mới $a_i$ và bảo toàn $K$ sau mỗi bước.
Các helper kiến tạo đóng sau chỉ dùng AST thô:
\begin{itemize}
  \item $setLiteralPlan(\bar a,\bar\tau,\upsilon)$ unwind danh sách
  $[setOut(coerce_C^{\tau_i\hookrightarrow\upsilon}(a_i))]_i$ rồi trả về
  phân biệt một alias;
  \item $unionPlan(P,Q,\upsilon)$ đổi tên hai cột đã coercion về cùng alias,
  ghép bằng $UnionAll$ và bọc một phép trả về phân biệt;
  \item $intersectPlan(P,Q,\upsilon)$ giữ một dòng đã coercion của $P$ khi
  $ExistsExpr$ trên $Q$ tìm thấy witness thỏa $setEq^\upsilon$, rồi trả về
  phân biệt;
  \item $asSetPlan(P)$ chiếu lại cột xuất của $P$ bằng $setOut$ và
  \texttt{DISTINCT};
  \item $uniquePlan(P,x,E)$ vật chất hóa $E$ một lần cho mỗi dòng của $P$ và
  trả về phép so sánh giữa số dòng nguồn phân biệt với
  $COUNT(DISTINCT\ setOut(E))$. Hai count là aggregate không grouping nên kế
  hoạch sinh đúng một dòng kể cả khi nguồn rỗng.
\end{itemize}
Mỗi helper nhập chính xác $FV$ của các kế hoạch con, cấp phát alias mới và trả
về schema một cột; vì vậy chúng không mở rộng ngôn ngữ đích.

Với kế hoạch vô hướng $P$, các alias mới $recv,out$ và bộ dựng thân thực thể
$F$ kết thúc bằng một phép chiếu tên $out$, định nghĩa
\[
\begin{split}
bottomArm(recv,b,out)&=Seq([Where(isBottom(recv)),Return(b\ AS\ out)]),\\
entityArm(recv,F,out)&=Seq([Where(NOT\ isBottom(recv))]+F(recv,out)),\\
withEntityReceiver(P,K,b,F)&=
ScalarPlan(X.prelude+[Call(K\cup\{recv\},B)],Alias(out)),
\end{split}
\]
trong đó $X=materialize(P,recv,K\cup\{recv\})$ và
$B=UnionAll(bottomArm(recv,b,out),entityArm(recv,F,out))$.
$F_{attr}$ thực hiện một phép khớp thuộc tính optional chuẩn và trả về giá trị
đã lưu; $F_{kind}$ trả về biểu thức tồn tại quan hệ phù hợp; còn $F_{cast}$ trả
về $recv$ khi biểu thức đó đúng và null nếu không. Theo thứ tự clause, dòng
bottom bị loại trước khi bất kỳ pattern thực thể nào được đánh giá. Các guard
bù nhau, thân thực thể đơn và import tường minh thiết lập $WFScalar$ cùng bảo
toàn alias đang sống.

Các phương trình khai triển biểu thức đúng kiểu ở mức ngôn ngữ host. Cụ thể,
\[
\begin{array}{rcl}
Expand_E(CQSetRelation)&=&ScalarPlan([],\,relationExpr),\\
Expand_E(CQCount(R))&=&ScalarPlan([],\,CountExpr(closeSet(Expand_S(R)))),\\
Expand_E(CQEmpty(R))&=&ScalarPlan([],\,emptyExpr),\\
Expand_E(CQIsUnique(R,x,E))&=&uniquePlan(Expand_S(R),x,Expand_E(E)),\\
Expand_E(CQAttribute(P,a))&=&withEntityReceiver(Expand_E(P),K,null,F_{attr}),\\
Expand_E(CQKindOf(P,C))&=&withEntityReceiver(Expand_E(P),K,false,F_{kind}),\\
Expand_E(CQCast(P,C))&=&withEntityReceiver(Expand_E(P),K,null,F_{cast}).
\end{array}
\]
Với kết quả tập, các phương trình bổ sung là
\[
\begin{array}{rcl}
Expand_S(CQSetLiteral(\bar P,\bar\tau,\upsilon))
 &=&setLiteralPlan(materializeList(Expand_E(\bar P),K),\bar\tau,\upsilon),\\
Expand_S(CQUnion(P,Q,\upsilon))
 &=&unionPlan(Expand_S(P),Expand_S(Q),\upsilon),\\
Expand_S(CQIntersection(P,Q,\upsilon))
 &=&intersectPlan(Expand_S(P),Expand_S(Q),\upsilon),\\
Expand_S(CQAsSet(P))&=&asSetPlan(Expand_S(P)).
\end{array}
\]
Do đó không phương trình $CQExpr$ nào trả về $CExpr$ trần tại nơi đối miền yêu
cầu $ScalarPlan$.

Phương trình nguồn giao receiver tập cho $Expand_S$ và nâng receiver thực thể
thành một dòng có guard, qua đó cài đặt $\asSources$. Mỗi phương trình còn lại
chỉ dùng constructor thô, cấp phát tên mới/tham số tất định và kết quả đệ quy đã
hoàn tất. Truy cập thuộc tính tạo phép khớp \texttt{ObjectHasAttribute}
optional trong prelude vô hướng. Navigation tạo pattern quan hệ không giới hạn
có hướng và các vị từ tường minh về \texttt{type}, liên kết, role, qualifier.
Các nút tồn tại và lực lượng chứa subquery đóng đã khai triển đệ quy. Điều kiện
vô hướng và tập tạo các nhánh \texttt{UnionAll} có guard bù nhau, cùng alias
được xuất và import tường minh. Phương trình invariant sinh phép khớp context,
prelude vị từ, bộ lọc vi phạm và phép trả về định danh phân biệt. Vì vậy không
placeholder import/alias trong ngoặc góc, khóa thuộc tính nội suy hay fragment
nhánh chưa khai triển nào thuộc đối miền của $Expand$.

\paragraph{Các suy diễn trường hợp khó của TXT4.}
Các trường hợp không trực tiếp của quy nạp hiện thực hóa đồng thời là:
\begin{enumerate}[label=(D\arabic*)]
  \item Với navigation có qualifier, áp dụng C1 lặp lại cho thấy việc vật chất
  hóa qualifier từ trái sang phải bảo toàn bộ qualifier ngữ nghĩa có thứ tự.
  Phép tuần tự hóa đơn ánh và R6 xác định chính xác các đích khớp; C5 loại các
  dòng quan hệ trùng.
  \item Với $IncludesAll$, subquery quan hệ thuộc bên trong đúng chính xác khi
  phần tử phía phải tương quan thuộc tập hữu hạn phía trái. Do đó lượng từ tồn
  tại bị phủ định bên ngoài phát biểu rằng không có phản ví dụ phía phải. Bỏ
  phủ định bên trong cho chứng minh tính rời nhau tương tự đối với $ExcludesAll$.
  \item Với iterator trên tập vô hướng, một dòng nguồn thỏa
  $\absC(r(a))=\encodeValue(s)=s$. C1 mở rộng tương ứng dòng bằng
  $x\mapsto a$ mà không gọi $\node(s)$, nên giả thiết quy nạp cho thân áp dụng
  đồng nhất với phần tử vô hướng và đối tượng.
  \item Với If có giá trị tập, điều kiện được vật chất hóa một lần. Chính xác
  một trong $truth(cc)$ và phủ định của nó đúng, hai nhánh trả cùng schema, và
  CY9 chỉ đóng góp các dòng của nhánh được chọn; phép chiếu ngoài phân biệt khôi
  phục ngữ nghĩa tập hữu hạn.
  \item Với Collect sinh bottom, $setOut$ đổi null thành token khác null.
  BR1--BR7 làm phép trừu tượng hóa khôi phục bottom và bảo đảm distinct/count
  không loại nó hoặc gộp nó với giá trị khác bottom.
  \item Với Let, đổi tên alpha rồi thế tránh bắt giữ bảo toàn định kiểu và ngữ
  nghĩa. Khai triển đệ quy dừng theo độ sâu lồng Let ngoài giảm nghiêm ngặt.
  \item Với thuộc tính, kind-of và cast trên bottom, $withEntityReceiver$ chỉ
  bật nhánh bottom không có pattern. Với thực thể, chỉ nhánh thực thể chạy và
  R4 hoặc R3 áp dụng. C1 và lập luận nhánh đơn thiết lập $WFScalar$ cùng bảo
  toàn mọi alias đang sống.
  \item Với Set literal, các kế hoạch phần tử được vật chất hóa từ trái sang
  phải. M3a/G3 và $coerce_C$ đưa chúng về kiểu $setJoin$ duy nhất; $setOut$
  giữ bottom và BR2/BR6 làm phép chiếu phân biệt đúng bằng tập ngoại diên của
  các giá trị phần tử.
  \item Với union, CY9 nối hai tập dòng sau coercion và BR6 loại trùng, nên kết
  quả là hợp ngoại diên. Với intersection, một dòng trái được giữ đúng khi
  CY4 tìm thấy dòng phải thỏa $setEq^\upsilon$; BR1--BR4 làm điều này tương
  đương membership có kiểu trong giao hai tập.
  \item Với $asSet$, giả thiết quy nạp đã cho một tập ngoại diên; chiếu lại
  bằng $setOut$ và \texttt{DISTINCT} không thêm hoặc bỏ giá trị theo BR6, nên
  đây là phép đồng nhất.
  \item Với $isUnique$, truy vấn nguồn có đúng một đại diện cho mỗi phần tử
  nguồn sau phép đóng tập. Thân được đánh giá một lần trên mỗi dòng; BR7/CY5
  làm so sánh lực lượng nguồn với lực lượng các giá trị thân phân biệt đúng khi
  và chỉ khi không có hai phần tử nguồn khác nhau có ảnh bằng nhau.
\end{enumerate}
Các suy dẫn D1--D11 cùng các trường hợp cơ sở/hợp thành trực tiếp suy ra TXT4 từ
C1--C5, R3--R6, BR1--BR7 và các tiên đề Cypher đã chọn; TXT4 không phải một giả
thiết độc lập về tính đúng đắn compiler.

\begin{lemma}[TXT5 Tính đóng và toàn phần của AST thô]
Mọi invariant đã chuẩn hóa được chấp nhận có một kết quả query-model và một kết
quả $Expand_I$ đóng, tới tương đương alpha theo tên mới. Ngoài ra,
\[
T_{\mathrm{TEXT}}(nva)=renderRaw(Expand_I(BuildCQInvariant(nva))),
\]
trong đó $renderRaw$ là pretty-printer toàn phần có đặt ngoặc và không có lựa
chọn ngữ nghĩa. Đây là đặc tả được chứng nhận; prototype sinh text trực tiếp
cần lập luận phù hợp chứng minh tương đương alpha với phép khai triển đã chỉ
định. Với $Q$ được sinh, nghĩa vụ dialect bổ sung là
$parse_{Cypher}(renderRaw(Q))\equiv_{\alpha}Q$.
\end{lemma}
\begin{proof}
TXT1 cho constructor query-model duy nhất. Quy nạp cấu trúc đồng thời trên
$CQExpr/CQSource/CQSet$ chứng minh tính toàn phần và invariant phạm vi kế hoạch
vì mọi đối số đệ quy là nút con thực sự, mọi alias được nhập/sinh đều tường minh
và mới. Let dùng độ đo lồng nhau quy chuẩn đã chứng minh. Đóng một trong hai kế
hoạch cho truy vấn thô không có lỗ. Round-trip parser là tiền đề dialect đã chọn
được hoàn thành bằng kiểm thử parser trên truy vấn sinh, không phải hệ quả chỉ
từ pretty-printing.
\end{proof}

Pretty-printer là toàn phần theo các phương trình constructor: alias, tham số,
thuộc tính, danh sách, biểu thức một/hai ngôi, case, hàm và
$ExistsExpr/CountExpr$ được render đệ quy với dấu ngoặc; pattern render nhãn,
hướng, kiểu hữu hạn và thuộc tính có tham số; mỗi constructor clause render từ
khóa và các nút con. $Seq$ nối các clause bằng dòng mới, còn
\[
renderRaw(UnionAll(Q_1,Q_2))=
renderRaw(Q_1)\mathbin{+}\texttt{ UNION ALL }\mathbin{+}renderRaw(Q_2).
\]
Với call tương quan,
\[
renderRaw(Call(K,Q))=
\texttt{CALL \{}\,renderRaw(injectImports(K,Q))\,\texttt{\}},
\]
trong đó $injectImports$ thêm import \texttt{WITH} tường minh vào đầu $Seq$ và
phân phối nó vào cả hai nhánh của $UnionAll$. Các phương trình này bao phủ mọi
constructor thô. Round-trip parser vẫn là tiền đề dialect đã chọn thay vì hệ
quả chỉ từ pretty-printing.

Luật gốc invariant dịch $nva$ thành Boolean dưới
$\alpha_0=\{self\mapsto s\}$, khớp $s$ qua \texttt{ObjectInstanceOf} tới
\texttt{:Class \{classKey:\$pC\}}, trong đó $\sigma(p_C)=key(C)$, lọc bằng
$NOT\ truth(c)$ và trả về \texttt{DISTINCT s.use\_id}. Đây là định nghĩa duy
nhất của $T_{\mathrm{TEXT}}$ cho Định lý~6.

\section{Các Định lý 0--6}

\begin{theorem}[Định lý 0: Biểu diễn đối tượng--đồ thị]
Với mọi cặp hợp lệ $(MM,M)$, nếu $G=\PhiEnc(MM,M)$ thỏa hợp đồng mã hóa đồ
thị, thì $G$ bảo toàn mọi quan sát liên quan đến kiểm tra mà $\OCLval$ sử dụng:
định danh đối tượng, tương ứng đối tượng--nút, sự phù hợp lớp, giá trị thuộc
tính, liên kết nhị phân, kết quả navigation và kết quả $allInstances$. Hình
thức hóa:
\[
\liftNode(\ObsObj(MM,M))=\ObsGraph(MM,\PhiEnc(MM,M)).
\]
\end{theorem}

\paragraph{Các bổ đề cần dùng.}
R1 Tính đơn ánh của đối tượng; R2 Tính chính xác đối tượng--nút; R3 Bảo toàn
kiểu; R4 Bảo toàn thuộc tính; R5 Bảo toàn liên kết; R6 Bảo toàn navigation; R7
Bảo toàn allInstances.

\begin{proof}
Định nghĩa $\Psi_{val}$ bằng các accessor chuẩn: nút đối tượng và
\texttt{use\_id} cho đối tượng, cạnh \texttt{ObjectInstanceOf} cho kiểu,
$I^{graph}_{attr}$ cho thuộc tính, quan hệ \texttt{Link*} cùng exact
\texttt{associationKey}, role và qualifier có hướng cho liên kết, $nav_{graph}$ cho
navigation, và quan hệ thuộc kiểu cho $allInstances$.

R2 cho song ánh giữa đối tượng nguồn và nút đối tượng nhìn thấy bởi kiểm tra qua
$\RepG$, đồng thời định nghĩa $\node(o)$; R1 cho tính đơn ánh. Vì thế
$\Psi_{val}$ không đọc thiếu và cũng không đọc nút giả. R3 bảo toàn sự phù hợp
lớp, R4 bảo toàn tra cứu thuộc tính. R5 bảo toàn liên kết theo hai bao hàm: mọi
liên kết UML đều được xuất và mọi liên kết đồ thị nhìn thấy bởi kiểm tra đều
được sinh từ một liên kết UML. Từ đó R6 bảo toàn navigation nhờ metadata hướng
và role đã phân giải. R7 bảo toàn $allInstances$ qua kế thừa được vật chất hóa.
Do từng thành phần của $\liftNode(\ObsObj(MM,M))$ bằng thành phần tương ứng của
$\ObsGraph(MM,G)$, đẳng thức ngoại diên của bộ chứng minh tương đương quan sát.
\end{proof}

\paragraph{Giải thích.}
Định lý không nói rằng đồ thị lưu mọi tạo tác UML. Nó chỉ khẳng định các quan
sát cần cho kiểm tra invariant được bảo toàn.

\begin{theorem}[Định lý 1: Tính đúng đắn của binder]
Nếu $e\in\OCLval$, $ast=\mathrm{Parse}(e)$,
$decode_{val}(ast)=e$, $\Adm_{MM}(ast)$ và
$T_{\mathrm{BIND}}(ast,MM)=b\in\BoundOCLval(MM)$, thì với mọi $M$ hợp lệ và
các môi trường gắn thẻ tương ứng $\rho_S\sim_{SB}\rho_B$,
\[
\sem{e}{S}(M,\rho_S)\sim_{SB}\sem{b}{B}(M,\rho_B).
\]
thì theo cầu nối xóa thẻ, hai ngữ nghĩa công khai của chúng bằng nhau.
\end{theorem}

\paragraph{Các bổ đề cần dùng.}
M2 Tính đóng của biểu thức Bound con và quy nạp cấu trúc; M3 Tính duy nhất của
kiểu chuẩn; M3a Tính hàm của $setJoin$; M4 Tính đúng đắn của admission và binding;
B1 Tính đúng đắn của
phân giải tên; B2 Tính đúng đắn của gán kiểu; B3 Tính tất định của phân giải
navigation; B4 Bảo toàn phạm vi từ vựng và che khuất; SB0 Tương ứng toán tử độc lập.

\begin{proof}
M4 đặt mọi kết quả thành công vào miền định lý hữu hạn có kiểu, M2 biện minh quy
nạp và M3 cố định kiểu trong từng phương trình cục bộ. Chứng minh bằng quy nạp
trên suy dẫn binding thành công. Chứng nhận ngữ nghĩa Source/Bound ở trên là
phân hoạch trường hợp đầy đủ, nên mỗi luật binder được chấp nhận có một phương
trình mô phỏng cục bộ. Các trường hợp cơ sở \texttt{self}, biến và literal là
trực tiếp. Với lời gọi thuộc tính, giả thiết quy nạp bảo toàn receiver; B1 phân
biệt thuộc tính với role liên kết, B2 bảo đảm dùng đúng kiểu, và B3 cho liên
kết, role, hướng cùng kiểu qualifier duy nhất của navigation. Các toán tử
Boolean, so sánh, số học và kiểu bảo toàn toán hạng tương ứng theo quy nạp; SB0
nối các kết quả toán tử được định nghĩa độc lập. Với \texttt{if} và
\texttt{let}, SB0 bảo toàn lựa chọn nhánh và mở rộng môi trường tương ứng. Với
iterator, B4 cho binding từ vựng tương ứng và giả thiết quy nạp áp dụng cho thân
dưới các môi trường đã mở rộng. Với Set literal, giả thiết quy nạp liên hệ từng
phần tử, M3a chọn duy nhất kiểu kết quả và SB0 áp dụng cùng phép nâng trước
khi lấy tập ngoại diên. $xor$ dùng cùng hai giá trị chân lý kiểm tra ở hai phía.
$isUnique$ được bảo toàn vì quan hệ phần tử là song ánh theo giả thiết quy nạp,
còn đẳng thức ảnh được SB0 phản ánh. Với union/intersection, hai tập con được
liên hệ ngoại diên sau cùng phép nâng $setJoin$; $asSet$ là đồng nhất. Đây là
toàn bộ các luật binder mới ngoài các họ đã xét. F1 riêng biệt chứng minh mọi biểu thức định kiểu
giới hạn có hình chiếu AST được chấp nhận; khi đã giả thiết admission, F1 không
phải tiền đề của quy nạp này.
\end{proof}

\paragraph{Giải thích.}
Binder không thay đổi ý nghĩa kiểm tra; nó phân giải tên và gắn kiểu cùng tham
chiếu siêu mô hình.

\begin{theorem}[Định lý 2: Trừu tượng hóa Đại số Kiểm tra]
Nếu $MM;\Gamma\vdash b:\tau$, $T_{\mathrm{VA}}(b)=va$ và
$\rho\models_M\Gamma$, thì
\[
erase_B(\sem{b}{B}(M,tag_B(\rho)))=\sem{va}{obj}(M,\rho).
\]
Tương đương, theo ký hiệu công khai dẫn xuất,
$\sem{b}{Bound}(M,\rho)=\sem{va}{obj}(M,\rho)$.
\end{theorem}

\paragraph{Các bổ đề cần dùng.}
M2 Tính đóng của biểu thức Bound con và quy nạp cấu trúc; M3 Tính duy nhất của
kiểu chuẩn; A1 Bảo toàn kiểu từ Bound sang VA; A2 Mô phỏng cục bộ từ Bound sang
VA; A3 Bảo toàn tra cứu môi trường; B4 Bảo toàn phạm vi từ vựng và che khuất.

\begin{proof}
M2 biện minh quy nạp cấu trúc trên biểu thức Bound hữu hạn và M3 cố định kiểu
chuẩn ở mỗi constructor. Bound \texttt{self}, biến và literal ánh xạ trực tiếp
sang \texttt{Self}, \texttt{Var}, \texttt{Literal} của VA; A3 cho cùng kết quả
tra cứu môi trường. Với thuộc tính và navigation, metadata đã phân giải được
sao chép sang nút VA tương ứng. Với Boolean, so sánh, số học, toán tử kiểu,
\texttt{if} và \texttt{let}, A2 cho cùng ngữ nghĩa toán tử phía đối tượng mà
không dùng mã hóa đồ thị. Với iterator và toán tử collection, tập nguồn được
bảo toàn theo quy nạp và thân được đánh giá dưới cùng môi trường mở rộng nhờ A3
và B4. Set literal dùng cùng phép nâng phần tử có kiểu trước khi VA dựng tập;
union/intersection dùng hợp/giao sau nâng, còn $asSet$ là đồng nhất. Với
$isUnique$, A3 cho cùng ảnh thân trên cùng tập nguồn và A2 phản ánh đẳng thức
của các giá trị ảnh, nên tính đơn ánh không đổi. Trường hợp $xor$ dùng phương
trình Boolean cục bộ của A2. A1 bảo đảm định kiểu đúng xuyên suốt.
\end{proof}

\paragraph{Giải thích.}
VA là một tầng trừu tượng: nó loại bỏ cú pháp riêng của parser nhưng giữ nguyên
ngữ nghĩa kiểm tra.

\begin{theorem}[Định lý 3: Bảo toàn chuẩn hóa]
Với mọi biểu thức VA định kiểu đúng $va$, nếu $nva=T_{\mathrm{NORM}}(va)$ thì:
\begin{enumerate}[label=(\roman*)]
  \item chuẩn hóa dừng;
  \item $nva$ ở dạng chuẩn;
  \item với diễn giải đối tượng và đồ thị $I$,
  \[
  \sem{va}{I}(\rho)=\sem{nva}{I}(\rho).
  \]
\end{enumerate}
\end{theorem}

\paragraph{Các bổ đề cần dùng.}
N1 Bảo toàn ngữ nghĩa của phép viết lại; N2 Bảo toàn kiểu của phép viết lại;
N3 Tính dừng của chuẩn hóa; N4 Bao phủ dạng chuẩn.

\begin{proof}
Bảo toàn ngữ nghĩa được chứng minh bằng phân trường hợp trên các luật viết lại.
Chẳng hạn, $Size(S)\to Count(S)$ giữ cùng lực lượng hữu hạn.
$ForAll(S,x,P)$ được viết lại thành $Not(Exists(S,x,Not(P)))$; hai vế tương
đương trên tập hữu hạn vì chân lý phổ quát chính là không tồn tại phản ví dụ.
$Reject(S,x,P)$ thành $Select(S,x,Not(P))$, giữ đúng các phần tử có vị từ không
validation-true. $Implies(A,B)$ thành $Or(Not(A),B)$ theo chính sách Boolean ở
mức kiểm tra. $Xor(A,B)$ thành
$Or(And(A,Not(B)),And(Not(A),B))$, đúng khi chính xác một toán hạng
validation-true. $Includes(S,y)$ thành kiểm tra tồn tại với đẳng thức có kiểu trên
tập hữu hạn; G3 bao phủ định danh thực thể và đẳng thức vô hướng tương thích.
$Excludes$ là phủ định trực tiếp của lượng từ tồn tại đó. Hai so sánh lực lượng
$Count(Nav(\cdots))>0$ và $Count(Nav(\cdots))=0$ lần lượt thành tồn tại và
không tồn tại trên kết quả navigation hữu hạn. N2 bảo toàn kiểu từng luật. N3
theo từ độ đo từ điển giảm trên mẫu nguồn, còn N4 theo từ phép duyệt dưới-lên
tất định.
\end{proof}

\paragraph{Giải thích.}
Chuẩn hóa thay đổi hình dạng biểu thức để đơn giản hóa việc sinh Cypher; nó
không thay đổi ý nghĩa kiểm tra.

\begin{theorem}[Định lý 4: Bảo toàn kiểm tra]
Cho $MM;\Gamma\vdash nva:\tau$. Nếu $G=\PhiEnc(MM,M)$ thỏa Định lý~0,
$\rho\models_M\Gamma$, $\eta\models_G\Gamma$ và
$\rho\sim_{\Phi}\eta$, thì
\[
\encodeValue(\sem{nva}{obj}(M,\rho))=\sem{nva}{graph}^{\tau}(G,\eta).
\]
Với các vị từ invariant:
\[
\boolval(\sem{nva}{obj}(M,self\mapsto o))
=
\boolval(\sem{nva}{graph}^{Boolean}(G,self\mapsto \node(o))).
\]
\end{theorem}

\paragraph{Các bổ đề cần dùng.}
Định lý 0; G1 Bảo toàn mã hóa môi trường; G2 Tính đơn ánh của mã hóa giá trị và
đồng cấu tập hữu hạn; G3 Bảo toàn đẳng thức có kiểu và toán tử vô hướng; G4
Tương thích với chính sách kiểm tra.

\paragraph{Chứng nhận tính đầy đủ theo constructor.}
Quy nạp được thực hiện trên suy dẫn định kiểu VA đã chuẩn hóa. Các dòng sau
phân hoạch toàn bộ văn phạm; ``đã loại'' nghĩa là N4 làm trường hợp ngoài cùng
không thể đạt tới sau chuẩn hóa.

\begingroup\scriptsize
\begin{longtable}{p{0.22\textwidth}p{0.31\textwidth}p{0.30\textwidth}p{0.10\textwidth}}
\toprule Constructor & Giả thiết quy nạp & Bổ đề chính & Trạng thái\\ \midrule
Literal/Self/Var & không có & quan hệ môi trường, mã hóa & đạt tới\\
SetLiteral & các phần tử & G2, ảnh tập hữu hạn ngoại diên & đạt tới\\
Attribute & receiver & R4 & đạt tới\\
Navigation & receiver và qualifier & R6, G4 $\asSources$ & đạt tới\\
AllInstances & không có & R3, R7 & đạt tới\\
Not/And/Or & các toán hạng & G3, chân lý G4 & đạt tới\\
Xor & các toán hạng & phương trình Boolean & đã loại\\
Implies & các toán hạng & phương trình Boolean & đã loại\\
đẳng thức/thứ tự/số học & các toán hạng & G3, hợp đồng vô hướng & đạt tới\\
If (vô hướng hoặc tập) & điều kiện và nhánh được chọn & G4, giả thiết quy nạp nhánh & đạt tới\\
Let & initializer và thân mở rộng & G1 & đạt tới\\
Exists & nguồn và vị từ mở rộng & G1, G2, G4 & đạt tới\\
ForAll & nguồn và vị từ mở rộng & tương đương phản ví dụ & đã loại\\
Select & nguồn và vị từ & G1, G2 & đạt tới\\
Reject & nguồn và vị từ & G1, G2 & đã loại\\
Collect & nguồn và thân & G1, G2 ảnh tập hữu hạn & đạt tới\\
IsUnique & nguồn và thân chiếu & G1, G2 phản ánh đẳng thức & đạt tới\\
Includes/Excludes & tập và phần tử & quan hệ thuộc G2 & đã loại\\
IncludesAll/ExcludesAll & cả hai tập & tập con/rời nhau G2 & đạt tới\\
Union/Intersection/AsSet & nguồn và đối số tập & hợp/giao/đồng nhất G2 & đạt tới\\
Count & nguồn & lực lượng G2 & đạt tới\\
Size & nguồn & từ đồng nghĩa Count & đã loại\\
IsEmpty/NotEmpty & nguồn & tính rỗng G2 & đã loại\\
KindOf & receiver & R3 & đạt tới\\
Cast & receiver & R3, thất bại thành bottom & đạt tới\\
\bottomrule
\end{longtable}
\endgroup

\begin{proof}
Chứng minh bằng quy nạp cấu trúc trên suy dẫn định kiểu của $nva$. Literal,
biến và \texttt{self} theo từ tương ứng môi trường. Set literal dùng giả thiết
quy nạp cho từng phần tử và tính ngoại diên của G2. Với truy cập thuộc tính,
receiver bottom cho bottom ở cả hai diễn giải theo phép hoàn chỉnh nguyên thủy;
receiver thực thể dùng R4. Navigation dùng $\asSources$ và R6; danh sách
qualifier thô có thứ tự được đánh giá trước, còn tuần tự hóa danh sách đơn ánh
làm đẳng thức thô tương đương đẳng thức payload. $AllInstances$ dùng R7. Các
trường hợp Boolean, so sánh và số học dùng giả thiết quy nạp cùng G3, kể cả đẳng
thức thực thể; G4 làm chân lý kiểm tra, $\finiteSet$ và $\asSources$ giao hoán
với mã hóa.

Với $Exists(S,x,P)$, giả thiết quy nạp cho tương ứng giữa các phần tử của
$\finiteSet(\sem{S}{obj})$ và $\finiteSet(\sem{S}{graph})$ qua
$\encodeValue$. G1 bảo toàn quan hệ sau khi mở rộng môi trường iterator, và giả
thiết quy nạp trên $P$ bảo toàn chân lý kiểm tra. $ForAll$, $Select$, $Reject$
và $Collect$ dùng cùng lập luận mở rộng môi trường. Quan hệ thuộc, tập con, rời
nhau, hợp, giao, $asSet$, tính duy nhất của ảnh, lực lượng và tính rỗng theo từ
tính đơn ánh và đồng cấu tập hữu hạn của G2. Khi kiểu phần tử hai phía khác
nhau nhưng $setJoin$ xác định, M3a và phương trình nâng kiểu của G3 cho cùng tập
sau coercion; vì thế lập luận hợp, giao và Set literal vẫn áp dụng ở kiểu kết
quả chung. Với $isUnique$, G1 ghép cặp các phần tử nguồn và G3 phản ánh đẳng
thức ảnh thân, nên tồn tại một cặp va chạm ở phía đối tượng khi và chỉ khi tồn
tại cặp va chạm tương ứng ở phía đồ thị.
\texttt{If} chọn các nhánh tương ứng nhờ chân lý được bảo toàn, còn
\texttt{Let} theo G1. Kind-of trên bottom là false và cast trên bottom là bottom
ở cả hai diễn giải. Với receiver thực thể, kind-of và cast theo R3, kể cả thất
bại thành bottom; \texttt{oclIsTypeOf} không có trường hợp constructor được
chứng nhận.
\end{proof}

\paragraph{Giải thích.}
Diễn giải đối tượng và diễn giải đồ thị của VA đã chuẩn hóa trùng khớp vì mọi
quan sát nguyên thủy mà VA sử dụng đều được phép mã hóa đồ thị bảo toàn.

\begin{theorem}[Định lý 5: Hiện thực hóa Cypher tham chiếu dưới các giả thiết Cypher]
Trên profile $CYPHER5_{val}$ đã chọn và một phiên bản server vượt qua các probe
dialect bắt buộc, nếu $MM;\Gamma\vdash nva:\tau$ và renderer bắt đầu từ ánh xạ
alias đơn ánh cùng ánh xạ tham số đúng đắn, thì kết quả không phải tập có suy
dẫn query-model loại biểu thức duy nhất, với khai triển toàn phần
$ScalarPlan(L,c)$ cho mọi tập alias đang sống hợp lệ $K$; kết quả tập có suy dẫn
loại tập duy nhất với khai triển $SetPlan(Q,a)$. Theo TXT5, mỗi kế hoạch được
sinh có khai triển AST thô đóng và thỏa phán đoán hiện thực hóa tương ứng. Với
$MM;\{self:C\}\vdash nva:Boolean$ và truy vấn invariant
$q=T_{\mathrm{TEXT}}(nva)$,
\[
\id(o)\in \returnedIds(q,G)
\iff
classOf(o)\mathrel{conformsTo}C\ \land\
\boolval(\sem{nva}{graph}^{Boolean}(G,self\mapsto \node(o)))=\mathrm{false}.
\]
\end{theorem}

\paragraph{Các bổ đề cần dùng.}
C1 Bảo toàn alias và phép chiếu đang sống; C2 Tính đúng đắn của tham số; C3
Hiện thực hóa truy cập đồ thị nguyên thủy; C3a Receiver thực thể an toàn với
bottom; C4 Hợp thành vị từ và lượng từ tồn tại; C5 Hợp thành tập hữu hạn và lực
lượng; C6 Hiện thực hóa lớp bọc invariant; các siêu thuộc tính ngoại diên
TXT1--TXT5; cùng các tiên đề ngữ nghĩa ngoại diên cho mảnh $Cypher_{val}$ được
sinh: tra cứu tham số/thuộc tính, \texttt{MATCH}, \texttt{WHERE},
\texttt{EXISTS}, \texttt{COUNT}, \texttt{WITH}, \texttt{UNWIND},
\texttt{CALL} tương quan, \texttt{UNION ALL}, \texttt{RETURN}, \texttt{CASE},
\texttt{coalesce} và \texttt{DISTINCT}.

\begin{proof}
Chứng minh bằng quy nạp trên suy dẫn query-model định hướng cú pháp và khai
triển AST thô khớp constructor, duy trì phán đoán hiện thực hóa thích hợp. TXT5
bảo đảm mỗi trường hợp sinh truy vấn đóng không có placeholder văn bản. Biến và
\texttt{self} được hiện thực hóa bởi alias, C1 ngăn bắt giữ. Literal được hiện
thực hóa bởi tham số, C2 bảo toàn giá trị. Thuộc tính dùng accessor đồ thị chuẩn
qua $withEntityReceiver$: bottom trả null mà không tạo node pattern, còn thực
thể đi vào nhánh optional-match duy nhất và dùng R4. Kind-of và cast dùng cùng
lập luận phân nhánh với R3. Navigation được hiện thực hóa bằng pattern
\texttt{MATCH} có biến quan hệ không giới hạn, điều kiện
\texttt{type(r) STARTS WITH 'Link'} và các vị từ trên
exact \texttt{associationKey}, role, hướng cùng
\texttt{sourceQualifiers}/\texttt{targetQualifiers}; C3 và hợp đồng đồ
thị cho cùng tập với $nav_{graph}$, tới các bản sao được C5 loại bỏ.

\texttt{Exists} được hiện thực hóa bởi \texttt{EXISTS \{...\}}, đúng chính xác
khi subquery có witness ngữ nghĩa theo C4. Selection và rejection dùng CY3;
collect đánh giá thân dưới mở rộng C1 và chiếu $setOut$ phân biệt. Quan hệ thuộc,
tập con và rời nhau dùng các truy vấn witness hoặc phản ví dụ với
$setEq^\tau$; C4--C5 chứng minh các lượng từ này đúng ngoại diên. Lực lượng và
tính rỗng lần lượt dùng CY4--CY5 cùng BR6--BR7.

Với Set literal, giả thiết quy nạp hiện thực hóa từng phần tử; M3a/G3 và
$coerce_C$ đưa chúng tới kiểu kết quả chung, còn BR4/BR6 giữ bottom và loại
đúng các phần tử trùng. Union dùng CY9 rồi chiếu phân biệt; intersection dùng
CY4 để giữ một phần tử trái khi có witness phải bằng nó theo
$setEq^\upsilon$; $asSet$ chỉ chiếu lại tập đã hiện thực hóa. Với $isUnique$,
CY5 và BR7 so sánh số phần tử nguồn với số giá trị thân phân biệt, đúng chính
xác khi ảnh thân đơn ánh. Các lập luận này là các trường hợp D8--D11 của TXT4.

Các toán tử Boolean, so sánh, số học và kiểu dùng CY1, C3a và C4; $xor$ cùng
các dạng Boolean nguồn đã bị chuẩn hóa được hiện thực hóa bởi các constructor
Boolean còn lại. Chân lý kiểm tra được cài đặt bằng
\texttt{coalesce(predicate,false)}.

Theo $Expand_I$ và C6, lớp bọc liệt kê thể hiện ngữ cảnh bằng quan hệ thuộc kiểu
đã vật chất hóa, giữ các thể hiện có vị từ không validation-true và trả về
\texttt{DISTINCT self.use\_id}. Do đó tập định danh trả về chính là tập vi phạm
phía đồ thị. C3 có chủ ý chỉ áp dụng cho nguyên thủy, nên phân tích trường hợp
cấu trúc ở đây, không phải bản thân C3, thiết lập hiện thực hóa cho biểu thức VA
đã chuẩn hóa bất kỳ trong phạm vi. Cuối cùng, $renderRaw$ chỉ thay đổi cách
trình bày cụ thể dưới tiền đề round-trip dialect của TXT5.
\end{proof}

\paragraph{Giải thích.}
Định lý phụ thuộc có điều kiện vào mảnh Cypher đã chỉ định. Đây không phải chứng
minh cho thực thi Cypher tùy ý.

\begin{theorem}[Định lý 6: Tương đương kiểm tra đầu-cuối]
Cho \texttt{context C inv R: e} là invariant với $e\in\OCLval$ và AST sau parse
thuộc $\ASTval$. Giả sử các tiền đề phạm vi chứng minh, đặc biệt
$ScalarClosed(e,M)$ và $BottomSeparated(MM,M,e)$, cùng chuỗi thành công:
\[
\mathrm{ast}=\mathrm{Parse}(e),\quad
s=decode_{val}(ast)=e,\quad
b=T_{\mathrm{BIND}}(\mathrm{ast},MM)\in\BoundOCLval(MM),\quad
va=T_{\mathrm{VA}}(b),\quad
nva=T_{\mathrm{NORM}}(va),\quad
q=T_{\mathrm{TEXT}}(nva),
\]
và $G=\PhiEnc(MM,M)$. Khi đó:
\[
\Viol_{OCL_{val}}(e,C,M)
=
\{\,o\in \Obj(M)\mid \id(o)\in \returnedIds(q,G)\,\}.
\]
\end{theorem}

\paragraph{Các kết quả cần dùng.}
Định lý 0; Định lý 1; Định lý 2; Định lý 3; Định lý 4; Định lý 5; tính đơn ánh
của định danh đối tượng; ngữ nghĩa tập cho các định danh được trả về.

F1 thiết lập khả năng biểu diễn biểu thức định kiểu giới hạn bằng AST được chấp
nhận, nhưng không cần dùng khi đã giả thiết chuỗi vận hành thành công này.

\begin{proof}
Ta chứng minh hai bao hàm.

\emph{Từ trái sang phải.} Giả sử $o\in \Viol_{OCL_{val}}(e,C,M)$. Khi đó
$classOf(o)$ phù hợp với $C$ và
\[
\boolval(\sem{e}{OCL_{val}}(M,self\mapsto o))=\mathrm{false}.
\]
Theo Định lý~1, giá trị này bằng giá trị Bound OCL. Theo Định lý~2, nó bằng giá
trị VA phía đối tượng. Theo Định lý~3, nó bằng giá trị VA đã chuẩn hóa phía đối
tượng. Theo Định lý~4, cùng chân lý kiểm tra thu được khi đánh giá $nva$ trên
$G$ với $self\mapsto\node(o)$. Vì vậy vị từ phía đồ thị không validation-true.
Theo Định lý~0, sự phù hợp kiểu kéo theo $\node(o)$ là một nút ngữ cảnh được
truy vấn invariant xét. Theo Định lý~5,
$\id(o)\in\returnedIds(q,G)$.

\emph{Từ phải sang trái.} Giả sử $\id_0\in\returnedIds(q,G)$. Theo đẳng thức
hiện thực hóa invariant trong Định lý~5, tồn tại $o\in\Obj(M)$ sao cho
$\id(o)=\id_0$, $classOf(o)$ phù hợp với $C$, và
\[
\boolval(\sem{nva}{graph}^{Boolean}
  (G,self\mapsto\node(o)))=\mathrm{false}.
\]
Tính đơn ánh của $\id$ làm đối tượng này duy nhất. Theo Định lý~4, vị từ đã
chuẩn hóa phía đối tượng không validation-true. Áp dụng ngược thứ tự Định
lý~3, Định lý~2 và Định lý~1, vị từ OCL ban đầu không validation-true. Do đó
$o\in \Viol_{OCL_{val}}(e,C,M)$.

Hai bao hàm đều đúng, nên hai tập vi phạm bằng nhau.
\end{proof}

\paragraph{Giải thích.}
Tương đương đầu-cuối thu được bằng cách hợp thành các kết quả bảo toàn cục bộ:
binding, trừu tượng hóa VA, chuẩn hóa, diễn giải đồ thị và hiện thực hóa Cypher
đều bảo toàn quyết định vi phạm cho từng đối tượng ngữ cảnh.

\section{Ranh giới cơ giới hóa Lean có chọn lọc}

Artifact \texttt{verification/lean/Ocl2CypherProof.lean} được ghim
vào Lean 4.32.2, proof contract \texttt{PC-2026-07-22.3} và SHA-256 hiện
hành của registry. Lean kernel kiểm tra 26 định lý được registry đặt
tên; checker cấm \texttt{sorry}, \texttt{admit}, khai báo \texttt{axiom}
của project và \texttt{opaque}. Axiom audit ghi công khai \texttt{propext}
ở 13 chứng minh; hai chứng minh dependent-scope trong số đó còn dùng core
\texttt{Quot.sound}. Chứng chỉ semantics của 11
rewrite và định lý T6 pointwise không phụ thuộc axiom. Mutation gate phải bắt được registry
hash cũ, placeholder, project-axiom, thiếu định lý, lỗi kiểu thật sự và core
axiom không được cho phép.

Phạm vi đã cơ giới hóa gồm đại số tập dạng predicate-set, tính
đơn ánh của \texttt{encodeValue} trên miền giá trị phẳng
bottom/scalar/entity/tập hữu hạn, đủ 11 rewrite family sau phép chiếu
validation-truth, relative normal form/idempotence/lexicographic root decrease,
N2 được index nội tại bởi \texttt{Bool}/\texttt{Nat}/\texttt{Elem}/\texttt{Set}
với erasure/type inference, bảo toàn biên binder bằng biểu diễn de Bruijn,
lexical named-to-de-Bruijn semantic correspondence, soundness và bảo toàn
scoped semantics của đúng ba disjunct trong Java capture guard trên Boolean
binder kernel,
lifting cấu trúc Boolean từ primitive agreement, cùng
hai bao hàm và mệnh đề pointwise của lõi trừu tượng Định lý~6.

Đây là PO-18 \emph{partial}. Artifact chưa biểu diễn nested collection,
OCL subtype/coercion phụ thuộc metamodel, tương ứng ngữ nghĩa phổ quát giữa
mọi production Java \texttt{OclIr} constructor và kernel đã chọn, hay Java
$T_{OPT}$--formal $T_{NORM}$ refinement,
quy nạp typed \texttt{OCL\_val} đầy đủ cho Định
lý~4, refinement của Java production, ngữ nghĩa Neo4j/Cypher hay
\texttt{AdapterAdequate}. Vì vậy kernel PASS chỉ xác nhận các mệnh đề
Lean đã chọn suy ra từ premise hiển thị; nó không phải định lý
correctness đầu-cuối cho prototype.

PO-19 cũng không được đóng chỉ vì CI đã được cấu hình. Workflow phải chạy trên
một commit sạch với mọi input required được track, tái capture hai manifest
runtime schema-v2 với \texttt{gitDirtyAtCapture=false}, rồi sinh/upload report
khớp đúng commit và registry. Điều kiện này đã được thỏa tại commit
\texttt{7b7bdd25}: clean baseline \texttt{f0938519}, Java 311/311, Lean 32/32,
hai mutation gate 6/6 và proof/build artifact đều PASS/upload thành công; do đó
PO-19 được discharged cho provenance CI/artifact. Source/PDF bài báo nằm ngoài
contract máy theo PO-20 và tiếp tục bị \texttt{md/} ignore.

\section{Kết luận}

Báo cáo đã thiết lập một chuỗi chứng minh khép kín tương đối với các tiền đề đã
phát biểu. Định lý~0 nối mô hình UML với các quan sát graph; Định lý~1--3 bảo
toàn lần lượt binding, phép trừu tượng hóa VA và chuẩn hóa; Định lý~4 nối hai
diễn giải VA; Định lý~5 hiện thực hóa diễn giải graph bằng mảnh Cypher được
chọn; và Định lý~6 hợp thành các kết quả đó thành đẳng thức của hai tập vi
phạm. Các trường hợp Set literal, $xor$, $union$, $intersection$, $asSet$ và
$isUnique$ được bao phủ bằng typing, coercion có kiểu, hiện thực hóa và các suy
diễn D8--D11 riêng.

Tính đúng đắn thu được là tính đúng đắn có điều kiện của đặc tả hình thức. Để
áp dụng trực tiếp cho một implementation cụ thể, implementation đó còn phải
thỏa $AdapterAdequate$, direct-text/formal/parser bridge cho phạm vi claim,
profile vô hướng, tách biệt bottom và các probe của phiên bản Cypher đã chọn.
Vì vậy kết quả phân biệt rõ
phần đã được chứng minh toán học với phần phải được giải phóng bằng conformance
test và bằng chứng thực thi.

Bằng chứng implementation hiện đã được tăng cường bằng kiểm tra R1--R7 trên ba
fixture, 47-case USE--Neo4j differential/premise run, selected parser-AST
projection 52/52 và ma trận CY1--CY9 9/9 trên Neo4j Kernel 2026.06.0
Enterprise. Selected conformance hiện chạy 285 test: 283 pass và đúng hai
freshness guard từ chối manifest schema-v1/contract-.2; mutation score 20/20;
runtime đạt 47/47 với 28 mixed/19 profile-tautology cùng ba
bottom/scalar/receiver probe.
Các kiểm tra hữu hạn này không chứng minh $AdapterAdequate$ phổ quát, full AST
hay universal plan closure, hoặc toàn bộ ngữ nghĩa Neo4j. Lean artifact đã
đóng các kernel chọn lọc nêu trên; hướng tiếp theo là mở rộng từ named/de-Bruijn
Boolean binder kernel sang full OCL/production-IR correspondence và Java refinement,
rồi quy nạp typed Định lý~4, trong khi giữ
Neo4j/Cypher và adapter ở biên premise/evidence được công bố rõ.

\section{Đồng bộ các tài liệu chứng minh}

Nguồn chuẩn duy nhất là
\texttt{md/research/formal-theorems-and-proofs.md}; canonical JSON block trong
file này cố định phiên bản hợp
đồng \texttt{PC-2026-07-22.3}, registry schema 4, A1--A9, M1--M5 (gồm M3a), statement
kernel và dependency chính xác của PC-T0--PC-T6, SF-01--SF-29, PO-01--PO-24,
runtime profile, 47 constructor và mười thuật ngữ implementation. File
\texttt{verification/contract/proof-contract-registry.json} và các block có
marker trong tài liệu này là hình chiếu xác định từ nguồn formal đó.
Kiểm tra \texttt{md/research/check-proof-sync.ps1} từ chối mã băm lỗi thời,
premise/dependency thiếu hoặc thừa, evidence mất, constructor/admission drift,
vocabulary drift và claim prototype quá sớm. Bảy mutation case, gồm cả trường
hợp chỉ sửa registry dẫn xuất, bảo đảm các lỗi này thật sự bị bắt; CI còn bắt
buộc biên dịch LaTeX. Cơ chế này chỉ thiết lập
đồng bộ và conformance tài liệu; nó không thay thế bất kỳ nghĩa vụ chứng minh
toán học nào.

\section{Ranh giới phát biểu trong bài báo}

Nên dùng:
\begin{quote}
Truy vấn invariant Cypher được sinh bảo toàn có điều kiện tập vi phạm đối với
mảnh kiểm tra OCL theo tập hữu hạn được hỗ trợ, dưới hợp đồng mã hóa đồ thị và
các giả thiết về mảnh Cypher đã phát biểu.
\end{quote}

Không nên dùng:
\begin{quote}
Chúng tôi chứng minh tính đúng đắn của OCL-to-Cypher cho toàn bộ OCL.
\end{quote}

\end{document}
